\documentclass[a4paper,titlepage,onecolumn,oneside,12pt,nenglish,titlepage=firstiscover]{scrbook}
% Um während der Entwurfsphase weniger leere Seiten zu haben, kann man oben "twoside" durch "oneside" ersetzen.

\usepackage{todonotes} % disable this package in final thesis!

\usepackage[english]{babel} % language
\usepackage[utf8]{inputenc} % encoding
\usepackage{amsmath,amsfonts,amssymb,amsthm} % mathematical expression, symbols, fonts, and the whole theorem environment engine
\usepackage{xcolor,graphicx} % graphics and colorings
\usepackage{tabularx,booktabs,multirow} % all about tables: booktabs gives the standard formatting, multirow allows cells to span over multiple rows.
%\usepackage{scrextend} % super powerful tool that helps a bit with everything
\usepackage[margin=0pt]{subcaption} % specific form for captions
\usepackage[autostyle]{csquotes} % helper for quoting

\usepackage{rotating} % if you want to rotate any object

\usepackage[backend=bibtex,
  style=alphabetic, % alphabetic ordered list
	hyperref=true, 
	natbib=true, % makes use of natbib
	backref=true, % additional writes for each bib-entry, on which pages its cites
	backrefstyle=three, % this style replaces e.g. 11,12,13 by 11-13
	sortcites=true, 
	maxbibnames=50,
	maxcitenames=3, % cites at most three authors---if more, only first author followed by et al. is written
 	firstinits=true, % nur Initialen anzeigen und auch nur nach diesen sortieren
% 	refsection=part,
% 	refsegment=chapter,
	isbn=false,url=true,doi=false] 
	{biblatex} % for references
\DeclareDelimFormat{finalnamedelim}{\addspace\bibstring{and}\space}

\bibliography{Abschlussarbeit_0495834} % add the .bib file you use here

\setlength\parindent{0pt}

% please fill in the following data:
\newcommand{\mytitle}{Automatic Optimization of Level Parameters for RSHash}
\newcommand{\mysubtitleA}{RSHash is a fast, compressed hash table for storing DNA k-mers}
\newcommand{\mysubtitleB}{recently introduced by the Reinert Group.}
\newcommand{\myname}{Bjarne David Jehn}
\newcommand{\thesisType}{Bachelor's Thesis}

\usepackage[pdfdisplaydoctitle,menucolor=orange!40!black,filecolor=magenta!40!black,urlcolor=blue!40!black,linkcolor=red!40!black,citecolor=green!40!black,colorlinks]{hyperref}
\hypersetup{pdftitle={\mytitle}, pdfauthor={\myname}}

%Caption formatting
\usepackage{caption}
\usepackage[font=small,labelfont=bf]{caption}
\captionsetup[figure]{
	format=plain,
	singlelinecheck=false
}
\captionsetup[table]{
	format=plain,
	singlelinecheck=false
}


%=========== graphics, in particular tikz
\usepackage{tikz}
\usetikzlibrary{arrows,decorations.pathmorphing,decorations.pathreplacing,backgrounds,positioning,fit,matrix}
\usetikzlibrary{shapes,calc,arrows.meta,positioning}

%=========== tables from csv files (helpful for experiments)
\usepackage{pgfplots}
\usepackage{siunitx}
\usepackage{pgfplotstable}
\usepgfplotslibrary{groupplots, statistics}
\pgfplotsset{compat=1.18}
\usepackage{csvsimple}

%=========== Pseudocode
\usepackage{algorithm}
\usepackage{algpseudocode}
\algnewcommand\algorithmicinput{\textbf{Input:}}
\algnewcommand\algorithmicoutput{\textbf{Output:}}
\algnewcommand\Input{\item[\algorithmicinput]}
\algnewcommand\Output{\item[\algorithmicoutput]}
\floatname{algorithm}{Algorithm}

%=========== \cref is a powerfull command that automatically states the current environment (e.g. Theorem, Lemma, etc.)
\usepackage[sort&compress,nameinlink,noabbrev]{cleveref}

%\usepgfplotslibrary{external}
%\tikzexternalize

%===========theorem environment
\theoremstyle{plain} % typical theorem-style (bold header, italic font)
\newtheorem{theorem}{Theorem}[chapter]
\newtheorem{lemma}[theorem]{Lemma}
\newtheorem{corollary}[theorem]{Korollar}
\newtheorem{observation}[theorem]{Beobachtung}
\newtheorem{proposition}[theorem]{Proposition}
\newtheorem{rrule}{Reduktionsregel}[section]

\crefname{rrule}{Reduktionsregel}{Reduktionsregeln} % how cleverref (\cref) has to handle these environments
\Crefname{rrule}{RR}{RRs} % with \Cref us can use what you defined here---somtimes useful if you want to abbreviate environment when citing (e.g. in tables or figures)

\theoremstyle{definition} % typical definition-style (bold header, normal font)
\newtheorem{definition}[theorem]{Definition}

\theoremstyle{remark} % typical theorem-style (italic header, normal font)
\newtheorem{example}{Example}
\newtheorem*{remark}{Remark}
\newtheorem{reduction}{Reduction}
\theoremstyle{plain}

%===========problem definition evironment: gets three arguments (1) problem name #1 (2) input specification #2 (3) question specification #3
\newcommand{\problemdef}[3]{
  \begin{center}
    \begin{minipage}{0.95\textwidth}
      \noindent
      \textsc{#1}
      
      \vspace{2pt}
      \setlength{\tabcolsep}{3pt}
      \begin{tabularx}{\textwidth}{@{}lX@{}}
        \textbf{Eingabe:} 		& #2 \\
        \textbf{Frage:} 	& #3
      \end{tabularx}
    \end{minipage}
  \end{center}
}



%=========== Nützliche Abkürzungen
\newcommand{\NN}{\mathbb{N}} % natürliche Zahlen
\newcommand{\RR}{\mathbb{R}} % reelle Zahlen
\newcommand{\QQ}{\mathbb{Q}} % rationale Zahlen
\newcommand{\ZZ}{\mathbb{Z}} % ganze Zahlen

\newcommand{\bigO}{\mathcal{O}} % Landausymbol

\newcommand\abs[1]{\left|#1\right|} % Betragsfunktion
\newcommand\norm[2]{\left\lVert#1\right\rVert_{#2}} % norm
\newcommand\enorm[1]{\norm{#1}{2}} % euclidean norm

%=========== commands for minimizer distribution
\newcommand{\sequence}{\mathcal{S}}
\newcommand{\listOfSuperKmers}{\textbf{skmers}}
\newcommand{\maxThreshold}{T}

%=========== commands for MINLP
\newcommand{\minimizerLength}{m}
\newcommand{\threshold}{t}
\newcommand{\level}{l}
\newcommand{\minimizer}{x}
\newcommand{\kmer}{q}

\newcommand{\access}{acc}
\newcommand{\contains}{rank}
\newcommand{\rank}{rank}
\newcommand{\select}{sel}
\newcommand{\compare}{comp}

\newcommand{\timeForOperationOnStructure}[2]{T_{#1_{#2}}} %#1:operation - e.g. rank, #2:structure - e.g. R

\newcommand{\varAverageFrequency}{\textbf{f}_{\textbf{avg}}}
\newcommand{\averageFrequency}{f_{avg}}

\newcommand{\possibleMinimizerLengths}{\mathcal{M}}
\newcommand{\possibleThresholds}{\mathcal{T}}

\newcommand{\indicatorm}{\mu}
\newcommand{\indicatort}{\tau}

\newcommand{\varMinimizerLength}{\textbf{\minimizerLength}}
\newcommand{\varThreshold}{\textbf{\threshold}}


%========== Title Formating
\titlehead{
\begin{minipage}{0.8\textwidth}
\textbf{Technische Universität Berlin}\\
Faculty IV: Electrical Engineering and Computer Science\\
%\textbf{Freie Universität Berlin}\\
%Department of Mathematics and Computer Science\\
%Algorithmic Bioinformatics\\

\end{minipage}
\hfill
\begin{minipage}{0.16\textwidth}
\includegraphics[width=\textwidth]{./figures/tu-logo}
\end{minipage}
\vspace{3.5cm}
}

\title{\mytitle}
\subtitle{$\left(\begin{gathered}
		\text{\mysubtitleA}\\
		\text{\mysubtitleB}
	\end{gathered}\right)$\vspace*{3cm}}
%\subtitle{(\mysubtitle)\vspace*{3cm}}

\author{
%	\Large{\textbf{\thesisType{}}}\vspace{0.1cm}\\
	\textbf{\large{\myname}}\vspace{0.1cm}\\
	\large{Thesis submitted in fulfillment of the requirements for the degree}\\
	\large{\grqq Bachelor of Science\grqq (B. Sc.) in the field of Computer Science}
	\vspace*{3cm}\\
		\small{
			\begin{tabular}{rl}
			    First Reviewer:& Prof. Dr. Knut Reinert\\
			    \hspace{2cm}Second Reviewer:& Prof. Dr. Stephan Kreutzer\\
			    Supervisor:& Jonas Schulte-Mattler
			\end{tabular}
		}
	\vspace*{1cm}
}

\date{} % only month and year
\begin{document}

\selectlanguage{english}

\maketitle

\section*{Declaration of Authorship}
\noindent
I hereby declare that the thesis submitted is my own, unaided work, completed without any external help. Only
the sources and resources listed were used. All passages taken from the sources and aids used, either
unchanged or paraphrased, have been marked as such.\\

Where generative AI tools were used, I have indicated the product name, manufacturer, the software version
used, as well as the respective purpose (e.g. checking and improving language in the texts, systematic
research). I am fully responsible for the selection, adoption, and all results of the AI-generated output I use.\\

I have taken note of the Principles for Ensuring Good Research Practice at TU Berlin dated 15 February 2023.
\url{https://www.static.tu.berlin/fileadmin/www/10002457/K3-AMBl/Amtsblatt\_2023/Amtliches\_Mitteilungsblatt\_Nr.\_16\_vom\_30.05.2023.pdf}\\

I further declare that I have not submitted the thesis in the same or similar form to any other examination
authority.\\

\begin{tabular}[h]{lll} Berlin, & \underline{17.11.2026} & \underline{\hspace{7cm}}\\
 & \small{Date} & \small{Signature}
\end{tabular}
\normalsize

\newpage
\section*{Abstract}
\todo[inline]{Draft}
RSHash is a static, compressed data structure for $k$-mers that answers lookup queries using a leveled minimizer scheme with thresholds to cut off frequent minimizers. Building it requires specifying a minimizer lengths and a threshold for every level used. Choosing optimal parameters to build RSHash for a given genetic sequence and specification of optimality remains an open problem as both its space consumption and the lookup time depend on these parameters which affect the distribution of minimizers. To address this, we first define properties of these distributions which are specifically relevant to the problem and propose a mixed integer non linear program (MINLP) modeling the space consumption and lookup time of RSHash aiming to combine it with a solver to find optimal parameters. We evaluate the proposed MINLP for a selection of benchmark datasets and can show that the model is highly accurate regarding RSHashs space consumption but less so for the lookup time.\todo{und was ist mit den optimalen parametern?} Furthermore, the time needed for computing the distribution stats required to run the model grows exponentially with the number of considered levels which means a similar complexity to that of a brute-force approach. 
\section*{Zusammenfassung}
\todo[inline]{Passt so? => Ins Deutsche übersezten}


\newpage
\tableofcontents
\newpage

\chapter{Introduction} \label{chap:introduction}
Genetic strands are chains of the nucleobases adenine (A), cytosine (C), guanine (G), thymine (T) and uracil (U). Adenine, cytosine and guanine are present in both DNA and RNA whilst thymine only appears in DNA and uracil is specific to RNA. A genetic sequence can thus be represented as a string over an alphabet of length four, allowing a compact representation of two bits needed to store each base. Despite this, genetic databases are often large as sequences can consist of several billion bases e.g. approximately $3.1 \cdot 10^9$ bases for the human genomes sequences. \\
The rapid growth of available genetic data due to better and faster sequencing techniques has led to a need for efficient algorithms for processing genetic sequences.\\

\textbf{$K$-mers} have emerged as an essential component of many approaches to encounter this issue. The $k$-mers of a sequence are all its substrings of length $k$, a sequence $S$ with $|S| \geq k$ therefore contains exactly $|S|-k+1$ $k$-mers. For example, the sequence $CGACAT$ consists of the 3-mers $CGA,GAC,ACA,CAT$.\\
Uses for $k$-mers include $k$-mer counting, frequency analysis, sequence indexing, genome assembly, sequence comparison, taxonomic classification, phylogeny reconstruction, protein sequence searching and alignment as well the construction of $k$-mer databases \cite{moeckel_survey_2024}.\\

\textbf{Sequence comparison} is used to compare given sequence or set of sequences against large genetic reference datasets, for example, to classify what species a collected DNA sample belongs to or for analyzing whether a patients DNA contains indicators linked to certain diseases.\\
One approach to this is checking whether each $k$-mer in the sequence(s) of interest is present in the reference dataset, also referred to as a lookup of the queried $k$-mers. On the other hand, a so called locate query answers not only whether queried $k$-mers are present but also tells all positions of the queried $k$-mers in the reference dataset. These positions are referred to as their offsets.\\

\textbf{A trade-off between space consumption and lookup time} is an inherent problem for data structures designed to answer lookup or locate queries. This can be demonstrated by considering the two extremes for data structures enabling a lookup or locate of a queried $k$-mer $q$ in a sequence $S$ with $|S| \gg k$ that doesn't contain any duplicate $k$-mers\footnote{The example could be adapted to sequences containing duplicate $k$-mers. We choose not to as this would add unnecessary complexity.}.\\
The most space efficient representation of all $k$-mers in $S$ is the sequence itself, needing approximately two bits per $k$-mer (bpk) as $|S| \gg k$. A lookup/locate query of $q$ would require scanning through the whole sequence, comparing every $k$-mer with $q$ and returning whether $q$ appears in the sequence, including its offset for the locate case. Both queries would terminate upon either finding a match or reaching the sequences end, requiring $O(|S|)$ time.\\
On the other hand, a fast lookup (or locate) could be achieved by storing every occurring $k$-mer in a hash set (locate: hash table), enabling lookup (or locate)  queries in $O(1)$ time. For this approach, every $k$-mer would have to be stored in the hash set (table) leading to a space consumption of $2k$ bits per $k$-mer in the lookup case and  $2k+log_2(|S|)$ bits per $k$-mer to enable locate queries which require storing each $k$-mers offset in the hash table.\\
Since both a lookup time in $O(|S|)$ as well as $2k$ bits per $k$-mer are not viable for large genomes, there has been much work on data structures enabling queries faster than $O(|S|)$ whilst also maintaining a certain compression.\\

\textbf{Using so called minimizers} is a common approach to achieve this. A minimizer of length $\minimizerLength$ is the smallest $\minimizerLength$-mer within a $k$-mer, regarding a total order over all strings of length $\minimizerLength$. For example, using the lexicographical order to determine the minimizer of length $2$ in the $6$-mer $CGACAT$ results in $AC$ because regarding that order, this is the smallest $2$-mer in $CGACAT$.\\
Minimizers have several properties which render them useful for this domain. Firstly, they are sparsely distributed because consecutive $k$-mers are likely to share the same minimizer. \Citet{zheng_improved_2020} show that in a random string $S$ with $|S| \gg k > \minimizerLength$ each minimizer of length $\minimizerLength$ is expected to cover an average of $(k-\minimizerLength+2)/2$ $k$-mers. Consecutive $k$-mers sharing the same minimizer are also referred to as super-$k$-mers. Even in sequences containing only unique $k$-mers, consecutive $k$-mers can share the same minimizer as $\minimizerLength < k$.\\
A second interesting property is the skewness of the distribution of minimizers in regards to their frequencies. The frequency of a minimizer is the number of super-$k$-mers which have this minimizer. In a given sequence only a few minimizers have high frequencies whilst most have small frequencies or do not appear at all.\\

Both the sparsity and skewness of minimizers have been observed by several previous works \cite{schulte-mattler_rshash_2026, pibiri_sparse_2022, jain_weighted_2020, chikhi_representation_2014}. We briefly describe Pibiri's SSHash \cite{pibiri_sparse_2022} and Schulte-Mattler's and Reinert's RSHash \cite{schulte-mattler_rshash_2026} in \cref{chap:previousWork} and specify relevant concepts and notation in \cref{chap:notation}. In \cref{chap:sequenceStatistics} we describe properties of genetic sequences which are particularly useful for modeling the space consumption and lookup time of the RSHash algorithm in \cref{chap:minlp}. Finally the proposed models accuracy and suitability for finding optimal level parameters is evaluated in \cref{chap:experiments}.

\chapter{Previous Work} \label{chap:previousWork}
\textbf{The SSHash algorithm} (Sparse and Skew Hashing of $k$-mers) introduced by \Citet{pibiri_sparse_2022} exploits these two properties by first computing all minimizers appearing in a sequence and then using a minimal perfect hash function to store every super-$k$-mers offsets in an array \textit{Offsets} in the order induced by their minimizers hashes. The $n$-th element of a second array \textit{Sizes} stores the starting index for the offsets in \textit{Offsets} belonging to the $n$-th minimizer. The sparsity of minimizers mainly effects the data structures space consumption: By indexing super-$k$-mers, the data structure is compressed by a factor of $(k-\minimizerLength+2)/2$ compared to indexing individual $k$-mers. On the other hand, the skewness of minimizers affects the locate time: Querying a $k$-mer $q$ requires computing its minimizer and then searching for it in every super-$k$-mer with that minimizer. On average, there are fewer super-$k$-mers per minimizer for larger minimizer lengths as high frequencies become less likely.\\
Therefore both space consumption and lookup/locate time of SSHash depend on the minimizer length: Shorter minimizer lengths $\minimizerLength$ mean higher compression ($(k-\minimizerLength+2)/2$) whilst the lookup/locate time is smallest for long minimizers as those result in lower frequencies. This demonstrates the challenge that comes with choosing an optimal minimizer length.\\

\newcommand{\infrequent}{\textit{infrequent} }
\newcommand{\frequent}{\textit{frequent} }
\newcommand{\hashfunction}[1]{h_{#1}}
\textbf{The RSHash algorithm} (Rank-Select Hash) is a minimizer based fast, compressed hash set for storing DNA $k$-mers that answers lookup and can be extended for supporting locate \cite{schulte-mattler_rshash_2026}.\\
It follows a minimizer based approach like SSHash; however, instead of using minimal perfect hashing, RSHash relies on bitvectors supporting rank and select queries. Its main difference to previous methods is the introduction of a layered multi-minimizer scheme aiming to reduce the cost of high frequencies whilst at also maintaining large compression.\\
The data structure works with one, two or three so called levels and a hash set (or hash table for answering locate queries). In order to build RSHash to answer lookup queries one has to specify the number of levels as well as a pair of level parameters for each level $\level$: The minimizer length $\minimizerLength^\level$ and a threshold $\threshold^\level$.\\

\citet{schulte-mattler_rshash_2026} show that RSHash allows significantly faster queries opposed to SSHash whilst remaining similarly compact. However, this is dependent of both the chosen minimizer lengths which influence the sparsity and skewness of minimizer distributions on each level and of the chosen thresholds which affect both the lookup time on each level as well as the percentage of $k$-mers having to be stored in the hash set.\\

\textbf{Building RSHash}: The first level is build by first computing each $k$-mers minimizer of length $\minimizerLength^1$ and counting the number of times each minimizer occurs in the sequences of interest. The minimizer computation uses a hash function $\hashfunction{1}$ to determine a random total order of the minimizers. We call minimizers that appear at least once, but no more than the threshold $\threshold^1$ \infrequent on the level.\\ 
A bitvector $R^1$ of lengths $|\Sigma|^{\minimizerLength^1}$ is used to mark the appearance of each distinct \infrequent minimizers in the order induced by $\hashfunction{1}$. For example, when using the lexicographical order and minimizers of length $\minimizerLength^1=2$, the first bit in $R^1$ would be set to one if $AA$ is an \infrequent minimizer in the sequences.\\ 
Each \infrequent minimizers occurrences in the sequences are stored in the offset vector $O^1$ in the order induced by $\hashfunction{1}$. Let $f$ for example be the frequency of the smallest \infrequent minimizer, that is the number of occurrences of the smallest minimizer for which the bit in $R^1$ is set. The first $f$ entries in $O^1$ are then the offsets of that minimizer whilst the following offsets belong to the second smallest \infrequent minimizer and so forth.\\
The boundaries of each minimizers offsets in $O^1$ are marked in the bitvector $S^1$ of length $|O^1|+1$. Following up on the example above, the first bit of $S^1$ would be set to one, indicating the start of the offsets belonging to the smallest \infrequent minimizer, this would be followed by $f-1$ zeros until the $f$-th bit in $S^1$ is again set to one, indicating the start of the second smallest \infrequent minimizers offsets and so forth. The last bit in $S^1$ is set to one marking the end of the largest \infrequent minimizers offsets.\\
If specified, all super-$k$-mers with \frequent minimizers on the first level are handled on the second level. This is build similar to the first level, creating the vectors $R^2, O^2$ and $S^2$ but using a different hash function $\hashfunction{2}$ for the computation of minimizers. Computing the minimizers with a different hash function again leads to a skewed distribution of minimizers. $K$-mers with \frequent minimizers on the first level, can therefore have \infrequent ones on the second. This holds even when the specified minimizer length $\minimizerLength^2$ and threshold $\threshold^2$ are equal to those on the first.\\
Further levels would be build analogously and only those $k$-mers with \frequent minimizers on all specified levels are stored in a fast hash set\footnote{When build to answer locate queries, a hash table is used to store the offsets of each $k$-mer with \frequent minimizers on all specified levels.}, sacrificing some space in return for a fast lookup. This is viable as only a small fraction of $k$-mers in the sequence of interest are \frequent on all levels for a good choice of the level parameters.\\

\textbf{However, the choice of "good" level parameters for a specific genome remains an open problem.} We already noted that using short minimizers leads to a larger compression but at the same time results in longer query times. As outlined, this contrary effect of the minimizer length on the space consumption and query time is an inherent problem of minimizer based $k$-mer indexing methods.\\
Existing methods use a number of approaches to deal with this problem. For example, \citet{pibiri_sparse_2022} chooses a relatively large minimizer length, accepting a smaller compression in order to maintain reasonable locate times.\\
The leveled approach introduced by \citet{schulte-mattler_rshash_2026} on the other hand aims to achieve high compression by using relatively short minimizers on lower levels whilst simultaneously mitigating the burden of \frequent minimizers by using thresholds to cut them off, bounding query times on every level. As outlined, the number of $k$-mers reduces with for subsequent levels and only a small percentage of the sequences $k$-mers end up being stored in the spatially expensive hash set. Therefore, choosing "good" level parameters to build the data structure for a given set of sequences currently requires an understanding of the RSHash algorithm and manual testing.

\chapter{Preliminaries}\label{chap:notation}
\newcommand{\numberOfLevels}{{L}}
\newcommand{\setOfLevels}{\mathcal{L}}

\newcommand{\levelParameters}[1]{p^{#1}}
\newcommand{\varLevelParameters}[1]{\textbf{p}^{#1}}
\newcommand{\parameterAccess}[3]{\levelParameters{#1}[#2^{#3}]}

\newcommand{\frequencyOf}[1]{freq(#1)}
\newcommand{\minimizerOf}[2]{min^{\hashfunction{#2}}_{\minimizerLength^{#2}}(#1)}

\newcommand{\sequences}[1]{S_{\levelParameters{#1}}}
\newcommand{\substring}[3]{#1[#2:#3)}
A sequence is a string over the DNA alphabet $\Sigma = \{A,C,G,T\}$. Let $i,j \in \NN$ with $i < j$ be zero based indices of a string $S$ of length $|S| \geq j$. We use $\substring{S}{i}{j}$ to denote the substring from the $i$-th to the $(j-1)$-th character of $S$.
The $k$-mers of a string $S$ with $|S|\geq k$ are all its substrings $S[i:i+k) \in \Sigma^k$ of length $k$ for $i \in \NN \wedge i \leq |S|-k$.\\
Let $\numberOfLevels \in \NN^+$ be the number of levels specified. We refer to the set of levels $\setOfLevels := \{ \level \in \NN^+ \mid \level \leq \numberOfLevels \}$. We use $\possibleMinimizerLengths \subseteq [1,k) \cap \NN$ and $\possibleThresholds \subset \NN^+$ to denote the sets of minimizer lengths and thresholds considered. We denote the Cartesian $n$-th power of a set $X$ with $X^n$ e.g. $X^3 = X \times X \times X$. A chosen minimizer length and threshold on a specified level $\level \in \setOfLevels$ are denoted $\minimizerLength^\level$ and $\threshold^\level$.\\
We use the tuple $\levelParameters{\level} := (\minimizerLength^1, ..., \minimizerLength^\level, \threshold^1, ..., \threshold^\level) \in \possibleMinimizerLengths^\level \times \possibleThresholds^\level$ to denote the level parameters up to level $\level$. For a specified tuple $\levelParameters{\level}$ of level parameters we use $\levelParameters{\level-1}$ to denote the level parameters $(\minimizerLength^1, ..., \minimizerLength^{\level-1}, \threshold^1, ..., \threshold^{\level-1})$ up to the previous level. Furthermore, we introduce the operator $\parameterAccess{\level}{\minimizerLength}{i} = \minimizerLength^i$ to denote retrieving the minimizer length specified for level $i$ from the level parameters and analogously $\parameterAccess{\level}{\threshold}{i}  = \threshold^i$.\\
We write $\sequences{\level-1} := \{S^\level_1,S^\level_2,...,S^\level_n\}$ to denote the set of sequences on level $\level$ in dependence of the level parameters $\levelParameters{\level-1}$ on the previous levels i.e. the $k$-mers which have \infrequent minimizers on all previous levels. The full set of input sequences when building RSHash do not depend on any parameters, are denoted with $\sequences{0}$ and are sometimes referred to as the input "text".\\
Let $\kmer \in \Sigma^{k}$ be a $k$-mer. The computation of its minimizer $\minimizerOf{\kmer}{\level}$ on level $l$ depends both on the specified minimizer length $\minimizerLength^\level$ as well as the hash function $\hashfunction{\level}$ on that level, which is not dependent of the level parameters but unique for every level. We denote the minimizers frequency with $\frequencyOf{\minimizerOf{\kmer}{\level}}$ which tells the combined number of the minimizer occurrences in a given set of input sequences.\\
We call the minimizer of $\kmer$ \infrequent on level $\level$ only if $1 \leq \frequencyOf{\minimizerOf{\kmer}{\level}} \leq \threshold^\level$ holds i.e. the minimizer occurs at least once but no more than the specified threshold on level $\level$. On the other hand, \frequent minimizers are those with frequencies larger than the threshold.\\
Let $B \in \{0,1\}^n$ be a one-indexed bitvector. For an index $i \in \NN^+$ with $i \leq n$, the vectors rank $\rank_B(i) = \sum_{j=1}^{i} B[j]$ tells the number of ones in $B$ up to index $i$. The select operation $\select_S(r) = min\{i \mid \rank_B(i) = r\}$ tells the index of the $r$-th $1$ in $B$.

\chapter{Sequence Statistics}\label{chap:sequenceStatistics}
As outlined in the \cref{chap:introduction}, the space consumption and query time of RSHash are highly dependent on the choice of level parameters as those influence the distributions of minimizers on all levels. We define properties of the minimizer distribution which will be needed for modeling the space consumption and lookup time in \cref{chap:minlp}. Specifically relevant for this is the behavior of minimizer distributions across several levels, which is unique to the approach proposed by \citet{schulte-mattler_rshash_2026} and therefore not considered in previous works.

\section{Minimizer Distribution} \label{sec:minimizerDistribution_severalLevels}
\newcommand{\distinctMinimizers}[1]{\mathcal{M}_{\levelParameters{#1-1}, \minimizerLength^{#1}}}
\newcommand{\distinctKmers}[1]{\mathcal{K}_{\levelParameters{#1-1}}}

\newcommand{\compactZV}[2]{#1[#2]}
\newcommand{\manualZV}[3]{#1[#2, #3]}
\newcommand{\ZVM}[1]{\manualZV{M}{\levelParameters{#1-1}}{\parameterAccess{#1}{\minimizerLength}{#1}}}
\newcommand{\ZVK}[1]{\manualZV{K}{\levelParameters{#1-1}}{\parameterAccess{#1}{\minimizerLength}{#1}}}

\newcommand{\probabilityDistribution}[1]{p_{#1}}

\newcommand{\CDF}[1]{F_{#1}}
\newcommand{\varCDF}[1]{\textbf{F}_{#1}}

Recall that we use $\levelParameters{\level} \in \possibleMinimizerLengths^\level \times \possibleThresholds^\level$ to describe the level parameters on all levels up to and including the level $\level$ and that $\sequences{0}$ is the full set of input sequences on the first level, which does not depend on any parameters. Also recall that the set $\sequences{\level-1}$ of input sequences on level $\level$ depends on the level parameters $\levelParameters{\level-1}$ on all previous level, as those affect, whether $k$-mers minimizer is \frequent or \infrequent.\\

Let $\distinctKmers{\level}$ be the set of distinct $k$-mers in a set of sequences $\sequences{\level-1}$. In other words, the set $\distinctKmers{\level}$ contains every distinct substring of length $k$ which appears in at least one of the sequences $S^\level_i \in \sequences{\level-1}$ which depend on the level parameters on previous levels:
\begin{equation}
	\distinctKmers{\level} := \{\kmer \in \Sigma^k \mid \exists S' \in \sequences{\level-1} ~ \exists (1 \leq i \leq |S'|-k) ~ S'[i:i+k) = \kmer \}
\end{equation}

Let $\distinctMinimizers{\level}$ be the set of distinct minimizers of lengths $\minimizerLength^\level$ on level $\level$ of the $k$-mers in $\distinctKmers{\level}$. The set of distinct minimizers is not only affected by the choice of level parameters on previous levels but also depends on the minimizer length on the current level:
\begin{equation}
	\distinctMinimizers{\level} := \{\minimizerOf{\kmer}{\level} \mid \forall \kmer \in \distinctKmers{\level}\}
\end{equation}
Let $\manualZV{M}{\levelParameters{\level-1}}{m}$ be a discrete random variable that describes the frequency of these minimizers:
\begin{align}
	\manualZV{M}{\levelParameters{\level-1}}{m} &: \distinctMinimizers{\level} \rightarrow \NN
	\\
	\manualZV{M}{\levelParameters{\level-1}}{m}(\minimizer) &= \frequencyOf{\minimizer}
\end{align}
We define its probability distribution $\probabilityDistribution{\manualZV{M}{\levelParameters{\level-1}}{\minimizerLength^\level}}$\footnote{This notation differs slightly from that used by \citet{schulte-mattler_rshash_2026} who do not introduce the probability distributions but instead use $p_M$ and $p_K$ to denote the cumulative distribution functions which we denote with $F_M$ and $F_K$}, which tells the probability of a minimizer having an exact frequency $f \in \NN$. The distribution of minimizers depends both on the level parameters $\levelParameters{\level-1}$ on previous levels, which affect the remaining $k$-mers on level $\l$ and the minimizer length $\minimizerLength^\level$ on the current level, which directly affects the computation of minimizers. We thus index the distributions with those parameters:
\begin{align}
	\probabilityDistribution{\manualZV{M}{\levelParameters{\level-1}}{\minimizerLength^\level}} &: \NN \rightarrow [0,1]
	\\
	\probabilityDistribution{\manualZV{M}{\levelParameters{\level-1}}{\minimizerLength^\level}}(f) &= \mathbb{P}(\manualZV{M}{\levelParameters{\level-1}}{\minimizerLength^\level} = f) = \frac{|\{\minimizer \in \distinctMinimizers{\level} \mid 
		\manualZV{M}{\levelParameters{\level-1}}{\minimizerLength^\level}(\minimizer)=f \}|}{|\distinctMinimizers{\level}|}
\end{align}
Using the random variables probability function we can derive its cumulative distribution function $\CDF{\manualZV{M}{\levelParameters{\level-1}}{m}}$, which returns the probability of a minimizer having a frequency that does not exceed an upper bound $f$: 
\begin{align}
	\CDF{\manualZV{M}{\levelParameters{\level-1}}{\minimizerLength^\level}} &: \NN \rightarrow [0,1]\\
	\CDF{\manualZV{M}{\levelParameters{\level-1}}{\minimizerLength^\level}}(f) &= \mathbb{P}(\manualZV{M}{\levelParameters{\level-1}}{\minimizerLength^\level} \leq f) = \sum_{i=1}^f \probabilityDistribution{\manualZV{M}{\levelParameters{\level-1}}{\minimizerLength^\level}}(i)
\end{align}

Whilst $\manualZV{M}{\levelParameters{\level-1}}{\minimizerLength^\level}$ and its distribution are useful for modeling values depending of a random (appearing) minimizer, RSHash is used to lookup $k$-mers, not minimizers. Therefore, we introduce the discrete random variable $\manualZV{K}{\levelParameters{\level-1}}{\minimizerLength^\level}$ which is defined for every distinct (appearing) $k$-mer on the level and describes the frequency each $k$-mers minimizer dependent of the level parameters $\levelParameters{\level-1}$ on all previous levels and the minimizer length $\minimizerLength^\level$ on the current one:
\begin{align}
	\manualZV{K}{\levelParameters{\level-1}}{\minimizerLength^\level} &: \distinctKmers{\level} \rightarrow \NN
	\\
	\manualZV{K}{\levelParameters{\level-1}}{\minimizerLength^\level}(\kmer) &= \frequencyOf{\minimizerOf{\kmer}{\level}}
\end{align}
Its probability distribution $\probabilityDistribution{\manualZV{K}{\levelParameters{\level-1}}{\minimizerLength^\level}}$ and its cumulative distribution function $\CDF{\manualZV{K}{\levelParameters{\level-1}}{\minimizerLength^\level}}$ are: 
\begin{align}
	\probabilityDistribution{\manualZV{K}{\levelParameters{\level-1}}{\minimizerLength^\level}} &: \NN \rightarrow [0,1]
	\\
	\probabilityDistribution{\manualZV{K}{\levelParameters{\level-1}}{\minimizerLength^\level}}(f) &= \mathbb{P}(\manualZV{K}{\levelParameters{\level-1}}{\minimizerLength^\level} = f) = \frac{|\{\kmer \in \distinctKmers{\level} \mid \manualZV{K}{\levelParameters{\level-1}}{\minimizerLength^\level}(\kmer)=f \}|}{|\distinctKmers{\level}|}
	\\
	\CDF{\manualZV{K}{\levelParameters{\level-1}}{\minimizerLength^\level}} &: \NN \rightarrow [0,1]\\
	\CDF{\manualZV{K}{\levelParameters{\level-1}}{\minimizerLength^\level}}(f) &= \mathbb{P}(\manualZV{K}{\levelParameters{\level-1}}{\minimizerLength^\level} \leq f) = \sum_{i=1}^f \probabilityDistribution{\manualZV{K}{\levelParameters{\level-1}}{\minimizerLength^\level}}(i)
\end{align}

Let $\levelParameters{\level}$ be a combination of level parameters on a level $\level$. Recall that we use $\parameterAccess{\level}{\minimizerLength}{\level}$ to denote retrieving the minimizer length specified for level $\level$ from the level parameters and analogously for $\parameterAccess{\level}{\threshold}{\level}$. For better readability we introduce a more compact notation to denote accessing a specific value of the probability distributions and their cumulative distribution functions in dependence of $\levelParameters{\level}$:
\begin{align}
	\probabilityDistribution{\compactZV{M}{\levelParameters{\level}}} &= \probabilityDistribution{\manualZV{M}{\levelParameters{\level-1}}{\parameterAccess{\level}{\minimizerLength}{\level}}}(\parameterAccess{\level}{\threshold}{\level})
	\\
	\probabilityDistribution{\compactZV{K}{\levelParameters{\level}}} &= \probabilityDistribution{\manualZV{K}{\levelParameters{\level-1}}{\parameterAccess{\level}{\minimizerLength}{\level}}}(\parameterAccess{\level}{\threshold}{\level})
	\\
	\CDF{\compactZV{M}{\levelParameters{\level}}} &= \CDF{\manualZV{M}{\levelParameters{\level-1}}{\parameterAccess{\level}{\minimizerLength}{\level}}}(\parameterAccess{\level}{\threshold}{\level})
	\\
	\CDF{\compactZV{K}{\levelParameters{\level}}} &= \CDF{\manualZV{K}{\levelParameters{\level-1}}{\parameterAccess{\level}{\minimizerLength}{\level}}}(\parameterAccess{\level}{\threshold}{\level})
\end{align}
For example, $\CDF{\compactZV{M}{\levelParameters{\level}}}$ denotes the probability of a minimizer being \infrequent on level $\level$ in dependence of the level parameters on every level as specified in $\levelParameters{\level}$.

\section{Average Frequencies} \label{sec:averagefrequencies}
Let $\Omega_{\manualZV{M}{\levelParameters{\level-1}}{\minimizerLength^\level}}$ be the sample space of the random variable $\manualZV{M}{\levelParameters{\level-1}}{\minimizerLength^\level}$ i.e. the set which contains all frequencies $f$ for which there is at least one minimizer $\minimizer$ on level $\level$ which occurs exactly $f$ times:
\begin{equation}
	\Omega_{\manualZV{M}{\levelParameters{\level-1}}{\minimizerLength^\level}} = \{f \mid \exists \minimizer \in \distinctMinimizers{\level} \frequencyOf{\minimizer} = f \}
\end{equation}
The average frequency of minimizers on a level $\level \in \setOfLevels$ is given by the expected value of the probability distribution of the random variable $\manualZV{M}{\levelParameters{\level-1}}{m}$:
\begin{equation}
	\mathbb{E}[\manualZV{M}{\levelParameters{\level-1}}{\minimizerLength^\level}] = \sum_{f \in \Omega_{\manualZV{M}{\levelParameters{\level-1}}{\minimizerLength^\level}}} \probabilityDistribution{\manualZV{M}{\levelParameters{\level-1}}{\minimizerLength^\level}}(f) * f
\end{equation}
For a given set of finite sequences the set of distinct minimizers on each level and for each combination of level parameters is finite by definition. It follows that for a given sequence and set of level parameters, the average frequency $\mathbb{E}[\manualZV{M}{\levelParameters{\level-1}}{\minimizerLength^\level}]$ can always be computed, because the sample space $\Omega_{\manualZV{M}{\levelParameters{\level-1}}{\minimizerLength^\level}}$ is then finite by definition.\\

More relevant to the optimization of level parameters is the average frequency of \infrequent minimizers i.e. the average frequency of minimizers whose frequency do not exceed the threshold $\threshold^\level$. We denote this value with $\averageFrequency(\levelParameters{\level})$ and can derive it from $\probabilityDistribution{\manualZV{M}{\levelParameters{\level-1}}{m}}$ by computing a weighted sum over all frequencies up until the threshold $\threshold^\level$ and then norming the result by dividing through the percentage of \infrequent minimizers:
\begin{equation} \label{equation:averageFrequency}
	\averageFrequency(\levelParameters{\level}) = \frac{1}{\CDF{\compactZV{M}{\levelParameters{\level}}}} \cdot
	\sum\limits_{f=1}^{\threshold^\level} \probabilityDistribution{\manualZV{M}{\levelParameters{\level-1}}{\parameterAccess{\level}{\minimizerLength}{\level}}}(f) \cdot f 
\end{equation}

\section{Numbers of Distinct \textit{Infrequent} Minimizers} \label{sec:nDIM}
\newcommand{\numberOfDistinctMinimizers}[1]{n_{DIM}(#1)}
\newcommand{\varNumberOfDistinctMinimizers}{\textbf{n}_{\textbf{DIM}}}
We define the number of distinct \infrequent minimizers $\numberOfDistinctMinimizers{\levelParameters{\level}}$ which also depend on all level parameters $\levelParameters{\level}$ of a level $\level$.  Their value can be computed py multiplying the total number of distinct minimizers on a level with the percentage of \infrequent ones:
\begin{equation}
	\numberOfDistinctMinimizers{\levelParameters{\level}} = |\distinctMinimizers{\level}| \cdot \CDF{\compactZV{M}{\levelParameters{\level}}}
\end{equation}

\section{Computing the Sequence Statistics} \label{sec:computingSequenceStatistics}
Computing the sequence statistics for a specified sequence $\sequences{0}$, set $\possibleMinimizerLengths$ of minimizer lengths and set $\possibleThresholds$ of thresholds to be considered requires the empirical probability distributions $\probabilityDistribution{\compactZV{M}{\levelParameters{\level}}}$ and $\probabilityDistribution{\compactZV{K}{\levelParameters{\level}}}$ for every combination of level parameters $\levelParameters{\level}$ on every specified level $\level$. One way of determining these is to use a recursion over the levels. Each recursive calls input consist of the current level $\level$, the number $\numberOfLevels$ of levels specified, the previous level parameters $\levelParameters{\level-1}$ and the set of input sequences in dependence of those parameters.\\
Upon being called, the recursive \Cref{alg:recursiveDistributionComputation} selects the first minimizer lengths $\minimizerLength \in \possibleMinimizerLengths$ in dependence of which the distribution stats for the input sequence on this level are computed. The function then loops over each threshold $\threshold \in \possibleThresholds$, removes all $k$-mers with \infrequent minimizers when using $(\minimizerLength, \threshold)$ from the input sequences and then calls the recursion for the next level, which receives the next level, the previous parameters extended by $(\minimizerLength, \threshold)$ and the modified sequences as input.
Only once this recursion is completed for every threshold $\threshold \in \possibleThresholds$, the next minimizer is picked from $\possibleMinimizerLengths$. Upon completion the cumulative distribution functions $\CDF{\compactZV{M}{\levelParameters{\level}}}$ and $\CDF{\compactZV{K}{\levelParameters{\level}}}$ (\cref{sec:minimizerDistribution_severalLevels}), the average frequencies $\averageFrequency(\levelParameters{\level})$ (\cref{sec:averagefrequencies}) and the numbers of distinct \infrequent minimizers $\numberOfDistinctMinimizers{\levelParameters{\level}}$ (\cref{sec:nDIM}) can be derived using the respective $\probabilityDistribution{\compactZV{M}{\levelParameters{\level}}}$ and $\probabilityDistribution{\compactZV{K}{\levelParameters{\level}}}$.\\
This approach is guaranteed to compute correct minimizer distributions for every set of level parameters on every level as it effectively performs the same steps as the RSHash algorithm. However, actually building RSHash for every combination of level parameters, would imply duplicate computation of the distributions on lower levels. This is not the case for the recursive approach, which considers every distribution exactly once.\\
Even though, this approach is therefore faster in practice, its runtime is still exponential in the number of levels specified. Possibilities of further speeding up this computation are outlined in \cref{chap:outlook}.
\begin{algorithm}[h]
	\begin{algorithmic} 
		\Input $
		\text{level } \level,
		\text{max Level } \numberOfLevels,
		\text{level parameters } \levelParameters{\level-1},
		\text{sequences } \sequences{\level-1}
		$
		\Function{recursiveMinimizerDistribution}{$\level, \numberOfLevels, \levelParameters{\level-1} \sequences{\level-1}$}
		\If{$\level = \numberOfLevels$} \State \textbf{return} \EndIf
		\For{$\minimizerLength \in \possibleMinimizerLengths$}
		\State $\probabilityDistribution{\manualZV{M}{\levelParameters{\level-1}}{\minimizerLength}},
		\probabilityDistribution{\manualZV{K}{\levelParameters{\level-1}}{\minimizerLength}} \gets getDistribution(\hashfunction{\level}, \sequences{\level-1}, m)$ \Comment $\hashfunction{\level}:$ hash function on level $\level$
		
		\For{$\threshold \in \possibleThresholds$}
		\State $\levelParameters{\level} \gets extend(\levelParameters{\level-1}, m,t)$	\Comment $extend((m^1,t^1), m^2,t^2) = (m^1,m^2,t^1,t^2)$
		\State $\sequences{\level} \gets getFrequentKmers(\level, \sequences{\level-1}, \minimizerLength, \threshold)$
		\State 
		\Call{recursiveMinimizerDistribution}{$\level+1, \numberOfLevels, \levelParameters{\level}, \sequences{\level}$}
		
		\EndFor
		\EndFor
		\EndFunction
	\end{algorithmic}
	\caption{Recursive Computation of Minimizer Distribution} \label{alg:recursiveDistributionComputation}
\end{algorithm}

\chapter{Mixed Integer Non Linear Program}\label{chap:minlp}
We model the optimization problem as a Mixed Integer Non Linear Programming (MINLP). In \cref{sec:minlp_oneLevel} we model the lookup time and space needed by RSHash on one level in dependence of the level parameters and minimizers distribution. We then extend this to the case of several levels in \cref{sec:minlp_severalLevels}.\\
In \cref{chap:experiments} we describe how to find the optimal level parameters for a given sequence using the MINLP.

\section{Model Input}\label{sec:minlp_modelInput}
\newcommand{\nkmers}{n_{kmers}}
The set of levels $\setOfLevels$, the sets of minimizer lengths $\possibleMinimizerLengths \subseteq [1,k) \cap \NN$ and thresholds $\possibleThresholds \subset \NN$ to consider have to be specified explicitly.
\\
Furthermore, the number $\nkmers$ of $k$-mers in the sequences the values $\CDF{\compactZV{M}{\levelParameters{\level}}}$, $\CDF{\compactZV{K}{\levelParameters{\level}}}$,  $\averageFrequency(\levelParameters{\level})$ and $\numberOfDistinctMinimizers{\levelParameters{\level}}$ for all possible combinations $\levelParameters{\level}$ of the level parameters on each level $\level \in \setOfLevels$ as defined in \cref{chap:sequenceStatistics} have to be passed as input to the MINLP. Their computation is further explained in \cref{chap:experiments}. 
\\
Lastly, a measure of the times 
$\timeForOperationOnStructure{\contains}{R}$,
$\timeForOperationOnStructure{\select}{S}$,
$\timeForOperationOnStructure{\access}{O}$,
$\timeForOperationOnStructure{\compare}{T}$,
$\timeForOperationOnStructure{check}{H}$
needed for a rank query on bitvector $R$, a select queries on bitvector $S$, accessing an offset in $O$, a comparison with the text $T$ at a random offset and checking whether a key exists in the hash set $H$. The computation of the minimizer distribution as well as the times needed to perform these actions on the specified data structures is outlined in \cref{chap:experiments}.

\section{Model for One Level}\label{sec:minlp_oneLevel}
Aiming to improve readability, we omit the superscript index indicating an instances belonging to a specific level as only one level is considered in this section e.g. the minimizer length is simply referred to as $\minimizerLength$ instead of $\minimizerLength^1$.
\\
In order to determine $\varMinimizerLength$ and $\varThreshold$ we use binary indicator variables $\indicatorm_\minimizerLength$ and $\indicatort_\threshold$ for each possible $\minimizerLength \in \possibleMinimizerLengths$ and $\threshold \in \possibleThresholds$:
\begin{align}
	\indicatorm_\minimizerLength &\in \{0,1\} \mid \forall \minimizerLength \in \possibleMinimizerLengths \label{variable:indicatorsM}\\
	\indicatort_\threshold &\in \{0,1\} \mid \forall \threshold \in \possibleThresholds \label{variable:indicatorsT}
\end{align} 

The indicator variables for each minimizer length should be mutually exclusive as $\varMinimizerLength$ can only take one distinct value. The same holds for the indicator variables for each threshold:
\newcommand{\mututallyExclusiveM}{\sum\limits_{\minimizerLength \in \possibleMinimizerLengths} \indicatorm_\minimizerLength &= 1}
\newcommand{\mututallyExclusiveT}{\sum\limits_{\threshold \in \possibleThresholds} \indicatort_\threshold &= 1}
\begin{align}
	\mututallyExclusiveM \label{constraint:mututallyExclusiveM}\\
	\mututallyExclusiveT \label{constraint:mututallyExclusiveT}
\end{align}

The variables $\varMinimizerLength \in \possibleMinimizerLengths$ and $\varThreshold \in \possibleThresholds$ are introduced to describe the optimal minimizer length and threshold.\\
There values can be modeled through a sum over the indicator variables multiplied with there corresponding values:
\newcommand{\optimalM}{\varMinimizerLength &= \sum\limits_{\minimizerLength \in \possibleMinimizerLengths} \indicatorm_\minimizerLength \cdot \minimizerLength}
\newcommand{\optimalT}{\varThreshold &= \sum\limits_{\threshold \in \possibleThresholds} \indicatort_\threshold \cdot \threshold}
\begin{align}
	\varMinimizerLength &\in \possibleMinimizerLengths \label{variable:optimalM}\\
	\varThreshold &\in \possibleThresholds \label{variable:optimalT}\\
	\optimalM \label{constraint:optimalM}\\
	\optimalT \label{constraint:optimalT}
\end{align}

We introduce the variable $\varCDF{M} \in [0,1]$ to take the value of $\CDF{\compactZV{M}{(\varMinimizerLength, \varThreshold)}}$ i.e. the percentage of minimizers of length $\varMinimizerLength$ in the input sequence, whose frequencies do not exceed $\varThreshold$ (\cref{sec:minimizerDistribution_severalLevels}).\\
It is not possible to directly evaluate this value as both $\varMinimizerLength$ and $\varThreshold$ are variables of the model. Instead we use the indicator variables to model a constraint forcing $\varCDF{M}$ to take the value of $\CDF{\compactZV{M}{(\minimizerLength, \threshold)}}$ for which $\indicatorm_\minimizerLength = 1$ and $\indicatort_\threshold = 1$:
\newcommand{\cdfMExpression}{
	\varCDF{M} &= \sum\limits_{(\minimizerLength, \threshold) \in \possibleMinimizerLengths \times \possibleThresholds} \CDF{\compactZV{M}{(\minimizerLength, \threshold)}} \cdot
	\indicatorm_\minimizerLength \indicatort_\threshold
}
\begin{align}
	\varCDF{M} &\in [0,1] \label{variable:percMinimizers}\\
	\cdfMExpression \label{constraint:percMinimizers}
\end{align}

Analogously we introduce $\varCDF{K} \in [0,1]$ to represent $\CDF{\compactZV{K}{(\varMinimizerLength, \varThreshold)}}$ and $\varAverageFrequency \in \RR^+_0$ to represent $\averageFrequency((\varMinimizerLength, \varThreshold))$ (\cref{sec:averagefrequencies}):
\newcommand{\cdfKExpression}{
	\varCDF{K} &= \sum\limits_{(\minimizerLength, \threshold) \in \possibleMinimizerLengths \times \possibleThresholds} \CDF{\compactZV{K}{(\minimizerLength, \threshold)}} \cdot
	\indicatorm_\minimizerLength \indicatort_\threshold
}
\newcommand{\varAverageFrequencyExpression}{
	\varAverageFrequency &= \sum\limits_{(\minimizerLength, \threshold) \in \possibleMinimizerLengths \times \possibleThresholds} \averageFrequency((\minimizerLength,\threshold)) \cdot
	\indicatorm_\minimizerLength \indicatort_\threshold
}
\begin{align}
	\varCDF{K} &\in [0,1] \label{variable:percKmers}\\
	\varAverageFrequency &\in \RR^+_0 \label{variable:averageFrequency}\\
	\cdfKExpression \label{constraint:percKmers}\\
	\varAverageFrequencyExpression \label{constraint:averageFrequency}
\end{align}

\subsection{Lookup Time}
In order to model the time needed for a lookup using RSHash, we closely follow the lookup procedure as described by \citet{schulte-mattler_rshash_2026}:\\
A $k$-mer $\kmer$ is looked up by first computing its minimizer $\minimizerOf{\kmer}{}$ and checking whether the bit $R[\minimizerOf{\kmer}{}]$ is set and if so simultaneously computing its rank $r = \rank_R(\minimizerOf{\kmer}{})$, taking $\timeForOperationOnStructure{\contains}{R}$ time (part of the model input \cref{sec:minlp_modelInput}).\\
For the percentage $\varCDF{K}$ (\cref{constraint:percKmers}) of the $k$-mers this bit is set and the lookup is completed by searching for $\kmer$ on the level. We introduce the variable $\timeForOperationOnStructure{level}{} \in \RR^+_0$ to describe the time needed for this.\\
If the bit $R[\minimizerOf{\kmer}{}]$ is not set, the lookup is completed through a lookup in the hash set $H$ taking $\timeForOperationOnStructure{check}{H}$ time (part of the model input \cref{sec:minlp_modelInput}). The variable $\timeForOperationOnStructure{lookup}{} \in \RR^+_0$ describing the lookup time can thus be modeled as:
\newcommand{\lookupTimeExpression}{
	\nonumber \timeForOperationOnStructure{lookup}{} =&~ \timeForOperationOnStructure{\contains}{R} + \varCDF{K} \cdot \timeForOperationOnStructure{level}{}\\
	&~\phantom{\timeForOperationOnStructure{\contains}{R}}+ (1-\varCDF{K}) \cdot \timeForOperationOnStructure{check}{H}}
\begin{align}
	\timeForOperationOnStructure{level}{} \in&~ \RR^+_0 \label{variable:levelTime}\\
	\timeForOperationOnStructure{lookup}{} \in&~ \RR^+_0 \label{variable:lookupTime}\\
	\lookupTimeExpression \label{constraint:lookupTime}
\end{align}

The search for those $k$-mers $\kmer$ for which the bit $R[\minimizerOf{\kmer}{}]$ is set continues on the level by using the minimizers previously computed rank  $r = \rank_R(\minimizerOf{\kmer}{})$ in the select queries $first = \select_S(r)$ and $last = \select_S(r+1)-1$. These tell the indices in Offsets $O$ where the minimizer’s first and last offset in the text $T$ are stored. The two select queries take $2\timeForOperationOnStructure{\select}{S}$ time.\\
Finally the presence of $\kmer$ in $T$ is checked at every minimizer offset in the interval $O[first:last]$ until $\kmer$ is either found or the end of the interval is reached. The average frequency $\varAverageFrequency$ of \infrequent minimizers is exactly the average length of this interval. In the lookup case, the comparison is terminated as soon as $\kmer$ is found for the first time, which is on average the case after $(\varAverageFrequency+1)/2$ comparisons with each comparison taking $\timeForOperationOnStructure{\compare}{T}$ time.\\
The following constraint models the complete time $\timeForOperationOnStructure{level}{}$ needed to search for a $k$-mers presence on the level:
\newcommand{\levelTime}{\timeForOperationOnStructure{level}{} &= 2\timeForOperationOnStructure{\select}{S} + (\varAverageFrequency+1)/2 \cdot (\timeForOperationOnStructure{\access}{O} + \timeForOperationOnStructure{\compare}{T})}
\begin{align}
	\levelTime \label{constraint:levelTime}
\end{align}
Note that the value $\timeForOperationOnStructure{lookup}{}$ is the average time for a positive RSHash lookup, that is the time to lookup a $k$-mer which was randomly sampled from the sequence. In practice, queried $k$-mers may or may not appear in the sequence - hence the need for a lookup. Looking up a $k$-mer which does not appear in the sequence (negative lookup) would always take $\timeForOperationOnStructure{\contains}{R} + \timeForOperationOnStructure{check}{H}$; however, we do not account for this as the ratio of positive and negative queries is impossible predict.\\

Without further constraints, solving the described MINLP for minimal $\timeForOperationOnStructure{lookup}{}$ results in a solution where the $\indicatorm_\minimizerLength$, $\indicatort_\threshold$ corresponding to the smallest $m \in \possibleMinimizerLengths$, $t \in \possibleThresholds$ will be set to one.\\
This is expected, as this leads to the smallest possible $\varCDF{K}$ which minimizes the $k$-mers having to be looked up in the text and maximizes the lookups in the (faster) hash set.\\
However, this results in maximal space consumption, as each entry in the hash set is significantly larger than an entry in $O$ as described in \cref{sec:problemDefinition}.\\
In order to account for that, we have to add constraints modeling the space consumption.

\subsection{Space Consumption}
Instead of modeling the total space used by RSHash we will focus on the bits per $k$-mer (bpk) used by the data structure as this metric allows to compare compression independent of the sequence length.\\
The space consumption of RSHash can be decomposed into the space needed for the data structures components: the text $T$\footnote{Recall that we use the 'text' $T$ to refer to the set of input sequences $\sequences{0}$ in order to avoid confusion with the bit vector $S$.}, the vectors $R$, $S$ and $O$ as well as the hash set $H$. We use one variable to describe the bits per $k$-mer needed by each of the components:
\begin{align}
	bpk_T &\in \RR_0^+ \label{variable:bpkT}\\
	bpk_R &\in \RR_0^+ \label{variable:bpkR}\\
	bpk_S &\in \RR_0^+ \label{variable:bpkS}\\
	bpk_O &\in \RR_0^+ \label{variable:bpkO}\\
	bpk_H &\in \RR_0^+ \label{variable:bpkH}
\end{align}

\textbf{The text $T$} is much longer than k. Because of this the bits per $k$-mer for storing the text $T$ can be set to a constant of $2$ bits needed to store one base:
\newcommand{\bpkT}{bpk_T &= 2}
\begin{align}
	\bpkT \label{constraint:bpkT}
\end{align} 
 
\textbf{Bitvector $R$} is a bitvector of length $4^\varMinimizerLength$ in which each bit set to $1$ marks the occurrence of one distinct \infrequent minimizer in the sequences. The amount of ones in $R$ can thus be computed by selecting the the With $\numberOfDistinctMinimizers{\levelParameters{(\minimizerLength, \threshold)}}$ being the number of distinct \infrequent minimizers in dependence of the specified level parameters, we can derive the value $\varNumberOfDistinctMinimizers$ for the chosen minimizer length $\varMinimizerLength$ and threshold $\varThreshold$:
\newcommand{\nDIMExpression}{
	\varNumberOfDistinctMinimizers &= \sum\limits_{(\minimizerLength, \threshold) \in \possibleMinimizerLengths \times \possibleThresholds}
	\numberOfDistinctMinimizers{(\minimizerLength, \threshold)} \cdot
	\indicatorm_\minimizerLength \indicatort_\threshold
}
\begin{align}
	\varNumberOfDistinctMinimizers &\in \RR^+_0 \label{variable:nDIM}\\
	\nDIMExpression \label{constraint:nDIM}
\end{align}
We note that for continuous sequences one could also model $\varNumberOfDistinctMinimizers$ by computing the number $\varCDF{K} \cdot \nkmers$ of $k$-mers with \infrequent minimizers and dividing it by the average frequency $\varAverageFrequency$ of minimizers multiplied with the ratio $(k-\varMinimizerLength+2)/2$ of $k$-mers covered per minimizer. This should decrease the models complexity as it would replace the non-linear constraint $\cref{constraint:nDIM}$ with a linear one. However, the ratio $(k-\varMinimizerLength+2)/2$ is only valid for continuous sequences $S$ of lengths $|S| \gg k > \varMinimizerLength$, rendering it useless both for sequences which have been preprocessed to remove duplicate $k$-mers and for \cref{sec:minlp_severalLevels} where the input on higher levels which consists of many relatively short super-$k$-mers.\\

Using Elias-Fano-Encoding, a bitvector of size $U$ with $n$ bits set to one can be compressed to
$n (\lceil log_2(U/n) \rceil + 3)$ \cite{okanohara2007practical}. Therefore, we can model the bits per $k$-mer needed to store $R$ by using this compression and dividing this by the total number $\nkmers$ of $k$-mers yields the constraint describing the bits per $k$-mer of $R$:
\newcommand{\bpkR}{bpk_{R} &= \frac{\varNumberOfDistinctMinimizers}{\nkmers} \cdot \left(\lceil 2\varMinimizerLength - log_2\varNumberOfDistinctMinimizers \rceil + 3 \right)}
\begin{align}
	\bpkR \label{constraint:bpkR}
\end{align}

\textbf{Bitvector $S$} has length $|O|+1 \approx |O|$ where every bit corresponds to the offset of one minimizer encountered in the sequences and one extra bit marks the end of all offsets - this last bit can be disregarded when modeling $bpk_S$ as its effect is insignificantly small.\\
Multiplying the number $\varNumberOfDistinctMinimizers$ of distinct \infrequent minimizers with their average frequency $\varAverageFrequency$ gives the length of the offset vector $O$. To model the bits per $k$-mer in $S$ we divide this by the number of $k$-mers in the sequence:
\newcommand{\bpkS}{bpk_{S}	&= \frac{\varNumberOfDistinctMinimizers \cdot \varAverageFrequency}{\nkmers}}
\begin{align}
	\bpkS \label{constraint:bpkS} 
\end{align}
Note that if one figured out, how to model the ratio between minimizers and $k$-mers on non-continuous sequences, $bpk_S$ could be modeled without the use of $\varNumberOfDistinctMinimizers$ by taking the percentage $\varCDF{K}$ of $k$-mers with infrequent minimizers and dividing it by that ratio.

\textbf{The offset vector $O$} contains one entry of size $log_2 |T|$ for the offset of each \infrequent minimizer.\\
Since $bpk_S$ already describes the number of offsets per $k$-mer, we can model the bits per $k$-mer in $O$ with:
\newcommand{\bpkO}{bpk_O &= bpk_S \cdot log_2 |T|}
\begin{align}
	\bpkO \label{constraint:bpkO}
\end{align}

\textbf{The hash set $H$} contains an entry $(key:\kmer)$ for each distinct $k$-mer $\kmer$ with an \infrequent minimizer, which are $1-\varCDF{K}$ of all $k$-mers in $T$.
Each entry has size $size(key) = 2k$ bit. We introduce the load factor\footnote{In practice, the load factor of a hash set is usually not known before building the set but rather a value used for re scaling. Therefore $\alpha=0.75$ is rather an estimate than the actual load factor of a possible hash set. We further asses the consequences of this in \cref{chap:experiments}.} $\alpha=0.75$.
The bits per $k$-mer required for the hash set can thus be modeled by:
\newcommand{\bpkH}{bpk_H &= (1-\varCDF{K}) \cdot \frac{2k}{\alpha}}
\begin{align}
	\bpkH \label{constraint:bpkH}
\end{align}

\textbf{The total bits per $k$-mer} $bpk_{total} \in \RR^+_0$ can be modeled by simply adding up the space requirements of each component:
\newcommand{\bpkTotal}{bpk_{total} &= bpk_T + bpk_R + bpk_S + bpk_O + bpk_H}
\begin{align}
	bpk_{total} &\in \RR^+_0 \label{variable:bpkTotal}\\
	\bpkTotal \label{constraint:bpkTotal}
\end{align}

\subsection{Optimization Goal}
\newcommand{\minimize}{minimize}
\newcommand{\boundTime}{Bound_{time}}
\newcommand{\boundSpace}{Bound_{space}}
Having defined all constraints regarding the lookup time and space consumption of RSHash, we can now construct a MINLP in order to find optimal level parameters.\\
We introduce three more external parameters to the model: The value $\minimize$ which specifies whether to minimize $\timeForOperationOnStructure{lookup}{}$ or $bpk_{total}$ and the upper bounds $\boundTime$, $\boundSpace$ for the lookup time and space consumption. 
\begin{align*}
	\minimize			&\in \{\timeForOperationOnStructure{lookup}{}, bpk_{total}\} \\
	\boundTime	&\in \RR_0^+\\
	\boundSpace		&\in \RR_0^+
\end{align*}
When using the MINLP, the upper bound belonging to the optimization goal should be set to $\infty$. In addition to the already defined ones, the following two constraints are added to the MINLP to enforce these bounds:
\begin{align}
	\timeForOperationOnStructure{lookup}{} &\leq \boundTime \label{constraint:boundTime}\\
	bpk_{total} &\leq \boundSpace \label{constraint:boundBpk}
\end{align}
One could for example search for level parameters minimizing the lookup time whilst not exceeding eight bits per $k$-mer by setting $\minimize \leftarrow \timeForOperationOnStructure{lookup}{}, \boundTime \leftarrow \infty, \boundSpace \leftarrow 8$. The constructed MINLP for this example is:
\begin{align*}
	INPUT: ~
	&\minimize = \timeForOperationOnStructure{lookup}{},
	\boundTime = \infty,
	\boundSpace = 8,
	\\
	&\setOfLevels = \{1\}, \possibleMinimizerLengths, \possibleThresholds, \nkmers
	\\
	&\CDF{\compactZV{M}{(\minimizerLength, \threshold)}}, \CDF{\compactZV{K}{(\minimizerLength, \threshold)}},
	\averageFrequency((\minimizerLength, \threshold)),
	\numberOfDistinctMinimizers{(\minimizerLength, \threshold)} ~ \forall (\minimizerLength, \threshold) \in \possibleMinimizerLengths \times \possibleThresholds,
	\\
	&\timeForOperationOnStructure{\contains}{R},
	\timeForOperationOnStructure{\select}{S},
	\timeForOperationOnStructure{\access}{O},
	\timeForOperationOnStructure{\compare}{T},
	\timeForOperationOnStructure{check}{H}
\end{align*}

\begin{align*}
	\text{MINIMIZE}\\
		&\minimize = \timeForOperationOnStructure{lookup}{}\\
	\text{SUBJECT TO}\\
	\mututallyExclusiveM &\Cref{constraint:mututallyExclusiveM} \\
	\mututallyExclusiveT &\Cref{constraint:mututallyExclusiveT} \\
	\optimalM &\Cref{constraint:optimalM}\\
	\optimalT &\Cref{constraint:optimalT}\\
	\cdfMExpression &\Cref{constraint:percMinimizers}\\
	\cdfKExpression &\Cref{constraint:percKmers}\\
	\varAverageFrequencyExpression &\Cref{constraint:averageFrequency}\\
	\levelTime &\Cref{constraint:levelTime}\\
	\lookupTimeExpression &\Cref{constraint:lookupTime}\\
	\timeForOperationOnStructure{lookup}{} &\leq \boundTime = \infty &\Cref{constraint:boundTime}\\
	\bpkT &\Cref{constraint:bpkT}\\
	\nDIMExpression &\Cref{constraint:nDIM}\\
	\bpkR &\Cref{constraint:bpkR}\\
	\bpkS &\Cref{constraint:bpkS}\\
	\bpkO &\Cref{constraint:bpkO}\\
	\bpkH &\Cref{constraint:bpkH}\\
	\bpkTotal &\Cref{constraint:bpkTotal}\\
	bpk_{total} &\leq \boundSpace = 8 &\Cref{constraint:boundBpk}\\
	%		\varMinimizerLength &= \sum\limits_{m \in \possibleMinimizerLengths} \indicatorm_\minimizerLength \cdot m \\
	%		\varThreshold &= \sum\limits_{t \in \possibleThresholds} \indicatort_\threshold \cdot t
\end{align*}

\begin{align*}
	\text{VARIABLES}\\
	\indicatorm_\minimizerLength &\in \{0,1\} \mid \forall m \in M &\Cref{variable:indicatorsM}\\
	\indicatort_\threshold &\in \{0,1\} \mid \forall t \in T &\Cref{variable:indicatorsT}\\
	\varMinimizerLength &\in \possibleMinimizerLengths &\Cref{variable:optimalM}\\
	\varThreshold &\in \possibleThresholds &\Cref{variable:optimalT}\\
	\varCDF{M} &\in [0,1] &\Cref{variable:percMinimizers}\\
	\varCDF{K} &\in [0,1] &\Cref{variable:percKmers}\\
	\varAverageFrequency &\in \RR_0^+ &\Cref{variable:averageFrequency}\\
	\timeForOperationOnStructure{level}{} &\in \RR^+_0 &\Cref{variable:levelTime}\\
	\timeForOperationOnStructure{lookup}{} &\in \RR_0^+ &\Cref{variable:lookupTime}\\
	bpk_T &\in \RR_0^+ &\Cref{variable:bpkT}\\
	bpk_R &\in \RR_0^+ &\Cref{variable:bpkR}\\
	bpk_S &\in \RR_0^+ &\Cref{variable:bpkS}\\
	bpk_O &\in \RR_0^+ &\Cref{variable:bpkO}\\
	bpk_H &\in \RR_0^+ &\Cref{variable:bpkH}\\
	\varNumberOfDistinctMinimizers &\in \RR^+_0 &\Cref{variable:nDIM}\\
	bpk_{total} &\in \RR_0^+ &\Cref{variable:bpkTotal}
\end{align*} 

\section{Model for Several Levels}\label{sec:minlp_severalLevels}
In the following, we adapt the constraints defined for using only one level in \cref{sec:minlp_oneLevel} to model the use of $\numberOfLevels \in \NN^+$ levels with $\setOfLevels = \{ \level \in \NN^+ \mid \level \leq \numberOfLevels \}$ being the set of levels used\footnote{We model the space consumption and lookup time for any positive natural number of levels even though RSHash currently only supports the use of one, two or three levels \cite{schulte-mattler_rshash_2026}.}.\\
When modeling several levels, we add a superscript $\level$ to variables in order to indicate there belonging to a level $\level \in \setOfLevels$ e.g. $\varMinimizerLength^1, \varMinimizerLength^2$ and $\varMinimizerLength^3$ denote the optimal minimizer lengths on level $1,2$ and $3$.\\

Unique sets of binary indicator variables $\indicatorm^\level_\minimizerLength$ and $\indicatort^\level$ are needed for the level parameters $\varMinimizerLength^\level$ and $\varThreshold^\level$ of each level $\level$.
The constraints regarding these variables are analog to those when using only one level:
\[ \begin{array}{rll}
	\indicatorm_\minimizerLength^\level &\in \{0,1\} & \mid \forall (\level,\minimizerLength) \in (\setOfLevels \times \possibleMinimizerLengths) \\
	\indicatort_\threshold^\level &\in \{0,1\} & \mid \forall (\level,\threshold) \in (\setOfLevels \times \possibleThresholds) \\
	\sum\limits_{m \in \possibleMinimizerLengths} \indicatorm_\minimizerLength^\level &= 1 & \mid \forall l \in \setOfLevels\\
	\sum\limits_{t \in \possibleThresholds} \indicatort_\threshold^\level &= 1 & \mid \forall l \in \setOfLevels\\
	\varMinimizerLength^\level &\in \RR^+_0 & \mid \forall l \in \setOfLevels\\
	\varThreshold^\level &\in \RR^+_0 & \mid \forall l \in \setOfLevels\\
	\varMinimizerLength^\level &= \sum\limits_{m \in \possibleMinimizerLengths} \indicatorm_\minimizerLength^\level \cdot \minimizerLength^\level & \mid \forall l \in \setOfLevels \\
	\varThreshold^\level &= \sum\limits_{t \in \possibleThresholds} \indicatort_\threshold^\level \cdot \threshold^\level & \mid \forall l \in \setOfLevels
\end{array} \]
For each level $l$ we model $\varCDF{M}^\level$, $\varCDF{K}^\level$ and $\varAverageFrequency^\level$ analogously to using a single level; however, it is necessary to sum over every possible combination of level parameters $\levelParameters{\level}$. The number of additions thus grows exponentially for an increasing number of levels.
\begin{align*}
	\varCDF{M}^\level &\in [0,1] &\mid \forall l \in \setOfLevels
	\\
	\varCDF{K}^\level &\in [0,1] &\mid \forall l \in \setOfLevels
	\\
	\varAverageFrequency^\level &\in \RR^+_0 &\mid \forall l \in \setOfLevels
	\\
	\varCDF{M}^\level &= \sum\limits_{\levelParameters{\level} \in \possibleMinimizerLengths^\level \times \possibleThresholds^\level}
	\CDF{\compactZV{M}{\levelParameters{\level}}} 
	\cdot
	\prod_{1\leq i \leq \level} \indicatorm^i_{\parameterAccess{\level}{\minimizerLength}{i}} \indicatort^i_{\parameterAccess{\level}{\threshold}{i}}
	& \mid \forall l \in \setOfLevels
	\\
	\varCDF{K}^\level &= \sum\limits_{\levelParameters{\level} \in \possibleMinimizerLengths^\level \times \possibleThresholds^\level}
	\CDF{\compactZV{K}{\levelParameters{\level}}}
	\cdot
	\prod_{1\leq i \leq \level} \indicatorm^i_{\parameterAccess{\level}{\minimizerLength}{i}} \indicatort^i_{\parameterAccess{\level}{\threshold}{i}}
	& \mid \forall l \in \setOfLevels
	\\
	\varAverageFrequency^\level &= \sum\limits_{\levelParameters{\level} \in \possibleMinimizerLengths^\level \times \possibleThresholds^\level}
	\averageFrequency(\levelParameters{\level})
	\cdot
	\prod_{1\leq i \leq \level} \indicatorm^i_{\parameterAccess{\level}{\minimizerLength}{i}} \indicatort^i_{\parameterAccess{\level}{\threshold}{i}}
	& \mid \forall l \in \setOfLevels
\end{align*}

\subsection{Lookup Time}
The search time on a level $l \in \setOfLevels$ can be modeled analogously to that on one level:
\begin{align*}
	\timeForOperationOnStructure{level}{}^\level & \in \RR^+_0 &\mid \forall l \in \setOfLevels\\
	\timeForOperationOnStructure{level}{}^\level &= 2\timeForOperationOnStructure{\select}{S} + \frac{\varAverageFrequency^\level+1}{2}  \cdot (\timeForOperationOnStructure{\access}{O} + \timeForOperationOnStructure{\compare}{T}) &\mid \forall l \in \setOfLevels
\end{align*}

For $\varCDF{K}^\level$ of the $k$-mers on a level $l \in \setOfLevels$ the lookup is completed by a search through the offsets on that level whilst the remaining $(1-\varCDF{K}^\level)$ of queries result in a lookup on the next level.\\
We can thus model the lookup time recursively:
\begin{align*}
	\timeForOperationOnStructure{lookup}{}^{l} &= 
	\begin{cases}
	\begin{aligned}
		\timeForOperationOnStructure{\contains}{R} &+ \varCDF{K}^\level \cdot \timeForOperationOnStructure{level}{}^\level\\
		&+ (1-\varCDF{K}^\level) \cdot \timeForOperationOnStructure{lookup}{}^{l+1}
	\end{aligned} &\mid l \in \setOfLevels\\
	\timeForOperationOnStructure{check}{H} &\mid l \notin \setOfLevels
\end{cases}
\end{align*}

Considering the set $\setOfLevels = \{ 1,2 \}$ of two levels as an example, this results in the following total lookup time:
\begin{align*}
	\timeForOperationOnStructure{lookup}{}^1 = \timeForOperationOnStructure{\contains}{R} +&~ \varCDF{K}^1 \cdot \timeForOperationOnStructure{level}{}^1\\
		+&~ (1-\varCDF{K}^1) \cdot \Big( \timeForOperationOnStructure{\contains}{R} + \varCDF{K}^2 \cdot \timeForOperationOnStructure{level}{}^2\\
		&~ \phantom{(1-\varCDF{K}^1) \cdot \Big( \timeForOperationOnStructure{\contains}{R^2}} + (1-\varCDF{K}^2) \cdot \timeForOperationOnStructure{check}{H} \Big)\\
\end{align*}
Inserting $\timeForOperationOnStructure{level}{}^1$ and $\timeForOperationOnStructure{level}{}^2$ yields:
\begin{align*}
	\timeForOperationOnStructure{lookup}{}^1 =
		\timeForOperationOnStructure{\contains}{R} +&~ \varCDF{K}^1 \cdot (2\timeForOperationOnStructure{\select}{S} + \frac{\varAverageFrequency^1+1}{2}  \cdot (\timeForOperationOnStructure{\access}{O} + \timeForOperationOnStructure{\compare}{T}))\\
			+&~ (1-\varCDF{K}^1) \cdot \Bigg(\timeForOperationOnStructure{\contains}{R} + \varCDF{K}^2 \cdot (2\timeForOperationOnStructure{\select}{S} + \frac{\varAverageFrequency^2+1}{2}  \cdot (\timeForOperationOnStructure{\access}{O} + \timeForOperationOnStructure{\compare}{T}))\\
			&~\phantom{(1-\varCDF{K}^1) \cdot \Bigg(\timeForOperationOnStructure{\contains}{R}} + (1-\varCDF{K}^2) \cdot \timeForOperationOnStructure{check}{H} \Bigg)
\end{align*}

\subsection{Space Consumption}
\textbf{The text $T$} is only being stored once, so the constant approximation $bpk_T \approx 2$ can be used for several levels as well.\\

\textbf{The number of distinct \infrequent minimizers} $\varNumberOfDistinctMinimizers^\level$ on each level $\level$ can modeled analogously to $\varCDF{M}^\level$, $\varCDF{K}^\level$ and $\varAverageFrequency^\level$:
\begin{align*}
	\varNumberOfDistinctMinimizers &\in \RR^+_0 &\mid \forall l \in \setOfLevels
	\\
	\varNumberOfDistinctMinimizers &= \sum\limits_{\levelParameters{\level} \in \possibleMinimizerLengths^\level \times \possibleThresholds^\level}
	\numberOfDistinctMinimizers{\levelParameters{\level}} 
	\cdot
	\prod_{1\leq i \leq \level} \indicatorm^i_{\parameterAccess{\level}{\minimizerLength}{i}} \indicatort^i_{\parameterAccess{\level}{\threshold}{i}}
	& \mid \forall l \in \setOfLevels
\end{align*}

\textbf{Vectors $R^\level$, $S^\level$ and $O^\level$:} The bits per $k$-mer needed to store the bit vectors $R^\level$ and $S^\level$ as well as the offset vector $O^\level$ on every level $\level$ can be modeled analogously to using only one level:
\begin{align*}
	bpk_{R^\level} & \in \RR^+_0 &\mid \forall l \in \setOfLevels
	\\
	bpk_{S^\level} & \in \RR^+_0 &\mid \forall l \in \setOfLevels
	\\
	bpk_{O^\level} & \in \RR^+_0 &\mid \forall l \in \setOfLevels
	\\
	bpk_{R}^\level &= \frac{\varNumberOfDistinctMinimizers^\level}{\nkmers} \cdot \left(\lceil 2\varMinimizerLength^\level - log_2\varNumberOfDistinctMinimizers^\level \rceil + 3 \right) &\mid \forall l \in \setOfLevels
	\\
	bpk_{S^\level} &= \frac{\varNumberOfDistinctMinimizers^\level \cdot \varAverageFrequency^\level}{\nkmers} &\mid \forall l \in \setOfLevels
	\\
	bpk_{O^\level} &= bpk_{S^\level} \cdot log_2 |T| &\mid \forall l \in \setOfLevels
\end{align*}

\textbf{The hash set $H$} only contains an entry for $k$-mers $\kmer$, which have frequent minimizers on all levels.
The bits per $k$-mer for the hash set can thus be modeled as 
\[ bpk_H = \left(\prod_{l=1}^{L} 1-\varCDF{K}^\level\right) \cdot \frac{64}{\alpha}\]\\

\textbf{The total bits per $k$-mer} for $L$ levels only differ to those on one level, as there is a vector $R,S$ and $O$ on each level:
\[ bpk_{total} = bpk_T + \sum_{l=1}^{L} \left(bpk_{R^\level} + bpk_{S^\level} + bpk_{O^\level}\right) + bpk_H\]

\chapter{Experiments}\label{chap:experiments}
In order to evaluate how suited the proposed MINLP (\cref{chap:minlp}) is for finding optimal level parameters, we use a set of reference datasets (\cref{sec:datasets}) and specify selection of considered level parameters in \cref{sec:consideredParameters}. In \cref{sec:computationOfModelInput} we show how to compute the model input and measure the actual space consumption and lookup time for every set of considered level parameters for each considered sequence. A selection of interesting observations regarding these measured values are pointed out in \cref{sec:observations}. We then use the computed input and the measured values in \cref{sec:modelAccuracy} to asses how accurate the modeled lookup time and space consumption are compared to the measured lookup time and space consumption. Finally we use the proposed model in \cref{sec:optimization} to find optimal level parameters using the SCIP solver \cite{AchterbergBertholdKochetal.2008, Achterberg2009, BestuzhevaChmielaMuellerSerranoVigerskeWegscheider2023}. Whether the resulting parameters are actually optimal in practice is highly dependent on how accurate the model is in the first place.

\section{Datasets} \label{sec:datasets}
We use the same datasets as \citet{pibiri_sparse_2022}. Those are the whole genomes of the atlantic cod (\textit{Gadus morhua}), the common kestrel (\textit{Falco tinnunculus}) and the whole GRCh38 human genome (\textit{Homo sapiens}) This choice allows a direct comparison with \citet{schulte-mattler_rshash_2026} who use the same datasets for benchmarking RSHash.\\
We assess the proposed MINLP both on the datasets themselves as well as on preprocessed versions without duplicate $k$-mers. For this a compacted de Bruijn Graph is built with BCALM2 \cite{chikhi_compacting_2016} whose maximal unitigs (non-branching paths) are extracted using the tool UST \cite{doi:10.1089/cmb.2020.0431}.\\


\section{Considered Parameters} \label{sec:consideredParameters}
For the following experiments we use $k$-mers of length $k = 31$ and the largest number of levels $\numberOfLevels = 3$ that is currently supported by RSHash.\\
RSHash currently works with thresholds up to $255$.
Technically, using the proposed MINLP to find optimal level parameters for a specified optimization goal, one would have to consider every minimizer lengths $1 \leq \minimizerLength < k$ and all thresholds $1 \leq \threshold \leq 255$. However, using the proposed model requires computing the distributions of minimizers for each possible combination of parameters. In practice using $3$ levels would require computing the distributions for $30^3 \cdot 255^3 \approx 448 \cdot 10^9$ combinations of level parameters which is is not viable.\\

Using minimizer lengths smaller than $15$ on any of the considered datasets leads to a large percentage of \frequent minimizers even when using the maximum threshold $\threshold=255$. As these result in a large hash set, minimizer lengths $\minimizerLength < 15$ can be disregarded when searching for level parameters that lead to relatively low query times whilst also keeping the data structure reasonably compact.
On the contrary, for minimizers of length $25$ the percentage of \infrequent minimizers is large for all considered genomes even when using a threshold as small as $1$. As the majority of appearing minimizers occurs only once, this percentage only changes insignificantly for thresholds $2 \leq t \leq 255$. Whilst this would mean a small hash set, the vector offset vector would get large as each minimizer only covers approximately $(31-25+2)/2 = 4$ $k$-mers.\\
For these reasons, considering minimizer lengths between fifteen and twenty five should be sufficient for analyzing the model. To further reduce the number of parameter combinations we only consider uneven minimizer lengths.
\begin{equation}
	\possibleMinimizerLengths = \{15,17,19,21,23,25\}
\end{equation}
Using exponentially increasing thresholds should be sufficient for analyzing the MINLP, because the cumulative distribution functions $\CDF{\compactZV{M}{\levelParameters{\level}}}$ and $\CDF{\compactZV{K}{\levelParameters{\level}}}$ have logarithmic growth. For each minimizer length $\minimizerLength \in \possibleMinimizerLengths$ the percentages of \infrequent minimizers changes only insignificantly for thresholds $128 \leq t \leq 255$ for each considered genome.
For an assessment of the proposed model it should therefore be sufficient to consider the thresholds $32,64$ and $128$ which are also those used by \citet{schulte-mattler_rshash_2026}: 
\begin{equation}
	\possibleThresholds = \{32,64,128\}
\end{equation}
With these sets of considered parameters, using $3$ levels would result in $6^3 \cdot 3^3 = 5,832$ combinations of level parameters. We reduce this even further by considering only non-decreasing parameter values on subsequent levels for both the minimizer lengths and thresholds. That means each combination of level parameters $\levelParameters{3} = (\minimizerLength^1, \minimizerLength^2, \minimizerLength^3, \threshold^1, \threshold^2, \threshold^3)$ fulfills $\minimizerLength^i \leq \minimizerLength^{i+1}$ and $\threshold^i \leq \threshold^{i+1}$ for $i \in \{1,2\}$. \citet{schulte-mattler_rshash_2026} argue that non-decreasing parameters for increasing levels should lead to high compression and a low expected maximum number of super-$k$-mers to verify on low levels.\\
For a non-decreasing sequence of length $k'$, chosen from $n$ values, the number of possibilities is $\binom{n+k'-1}{k'}$. That means when choosing three minimizer lengths out of six possible values and three thresholds out of three possible values, the total number of parameter combinations is $\binom{6+3-1}{3} \cdot \binom{3+3-1}{3} = 560$. Computing the minimizer distribution for each of those combinations on all considered genomes is viable.

\section{Computing the Model Input} \label{sec:computationOfModelInput}
In order to evaluate the accuracy of the proposed MINLP, we first show how to compute the model input for each of the raw/preprocessed datasets which we refer to as the set of sequences $\sequences{0}$. This is done for $\numberOfLevels = 3$ levels and non-decreasing level parameters as described in \cref{sec:consideredParameters}.\\
The required stats are the total number $\nkmers$ of $k$-mers in the input sequence and the cumulative distributions of minimizer frequencies $\CDF{\compactZV{M}{\levelParameters{\level}}}$ and $\CDF{\compactZV{K}{\levelParameters{\level}}}$, the average frequencies $\averageFrequency(\levelParameters{\level})$ and the number of distinct \infrequent minimizers $\numberOfDistinctMinimizers{\levelParameters{\level}}$ for every non-decreasing combination of level parameters $\levelParameters{\level} \in \possibleMinimizerLengths^\level \times \possibleThresholds^\level$ on every level $1 \leq \level \leq 3$.\\

These stats can be brute-forced for every dataset considered by iterating over each non-decreasing combination of the specified level parameters and actually building RSHash with those parameters. We use the RSHash algorithm implemented by \cite{schulte-mattler_rshash_2026} which is available at \url{https://github.com/jonsmcode/rshash}.\\ Whilst there are more efficient approaches to only compute the model input, we chose this approach as it not only computes an accurate model input but also allows to measure the space consumption of the data structures components, namely the bits per $k$-mer of the vectors $R, S$ and $O$ on each level as well as the bits per $k$-mer to store the text and hash set which are only stored once and not per level.\\
The average lookup time $\timeForOperationOnStructure{lookup}{}$ is measured by for $100000$ $k$-mers randomly selected from the sequence. Furthermore, we measure an average for each individual operation time $\timeForOperationOnStructure{\contains}{R}$,
$\timeForOperationOnStructure{\select}{S}$,
$\timeForOperationOnStructure{\access}{O}$,
$\timeForOperationOnStructure{\compare}{T}$ and
$\timeForOperationOnStructure{check}{H}$. \Cref{algo:computeMinDistr} shows the described algorithm.
\begin{algorithm}[h]
	\caption{Brute-force Computation of Model Input and Measurements Needed for Evaluation.} \label{algo:computeMinDistr}
	\begin{algorithmic} 
		\Input $
		\text{sequences } \sequences{0},
		\text{max Level } \numberOfLevels,
		\text{considered minimizer lengths and thresholds } \possibleMinimizerLengths, \possibleThresholds
		$
		\Function{statsComputation}{$\sequences{0}, \numberOfLevels, \possibleMinimizerLengths, \possibleThresholds$}
		\State $\nkmers \gets ccountKmers(\sequences{0})$
		\For{$\levelParameters{L} \in \possibleMinimizerLengths^L \times \possibleThresholds^L$}
		\If{$\levelParameters{L}.isDecreasing()$} \State \textbf{continue} \EndIf
		\State $RSHash \gets buildWithStats(\sequences{0}, \levelParameters{L})$
		\State $\CDF{\compactZV{M}{\levelParameters{\level}}},
		\CDF{\compactZV{K}{\levelParameters{\level}}},
		\averageFrequency(\levelParameters{\level}), \numberOfDistinctMinimizers{\levelParameters{\level}} \gets RSHash.getDistributionStats()$
		\State $spaceConsumptions [\levelParameters{L}] \gets RSHash.getSpaceConsumptions()$
		\State $\timeForOperationOnStructure{lookup}{} [\levelParameters{L}] \gets RSHash.benchLookup(noSamples=100000)$
		\State $operationTimes [\levelParameters{L}] \gets RSHash.getOperationTimes(noSamples=100000)$
		\EndFor
		\EndFunction
	\end{algorithmic}
\end{algorithm}

\section{Observations Regarding the Computed Stats} \label{sec:observations}
\subsection{Minimizer Distribution}
Recall that RSHash uses several levels aiming to repeatedly take advantage of the sparse and skewed distribution of minimizers. This requires using a different hash functions to compute minimizers on each level as using the same function on every level would mean that give equal minimizer length and threshold, those minimizers which are \frequent on a level would also be \frequent on the next.
We observe that despite of the minimizers being computed using a different hash function, their distributions on levels $\level > 1$ remain highly dependent on the distributions on all previous levels. \Cref{fig:multilevelDistribution} shows this for the raw genomes using the level parameters ($\levelParameters{\level} = (17,17,17,32,32,32)$) to highlight the observation. We observe a similar dependency for differing parameters but choose to show it for identical ones, as it highlights the percentage being smaller on every subsequent level.\\
\begin{figure}[h]
	\centering
	\begin{tikzpicture}
	\begin{groupplot}[
		group style={
			group size=3 by 1,
			horizontal sep=0.07\textwidth,
		},
		width=0.33\textwidth,
		height=5cm,
		ymin=0, ymax=1,
		xlabel={$f$},
		]
		
		% First panel
		\nextgroupplot[
			title={Kestrel (raw)},
			ylabel={$F_M$},
			legend to name=sharedlegend,
			legend columns=3,
			]
		\addplot+[only marks, mark=*, mark size=1.5pt] table[col sep=comma, x=frequency, y=F_M_1]
		{./data/F_M_Falco_tinnunculus.FalTin1.0.dna.toplevel.csv};
		\addlegendentry{Level $1$}
		\addplot+[only marks, mark=*, mark size=1.5pt] table[col sep=comma, x=frequency, y=F_M_2]
		{./data/F_M_Falco_tinnunculus.FalTin1.0.dna.toplevel.csv};
		\addlegendentry{Level $2$}
		\addplot+[only marks, mark=*, mark size=1.5pt] table[col sep=comma, x=frequency, y=F_M_3]
		{./data/F_M_Falco_tinnunculus.FalTin1.0.dna.toplevel.csv};
		\addlegendentry{Level $3$}
		
		% Second panel
		\nextgroupplot[title={Cod (raw)}]
		\addplot+[only marks, mark=*, mark size=1.5pt] table[col sep=comma, x=frequency, y=F_M_1]
		{./data/F_M_Gadus_morhua.gadMor3.0.dna.toplevel.csv};
		\addplot+[only marks, mark=*, mark size=1.5pt] table[col sep=comma, x=frequency, y=F_M_2]
		{./data/F_M_Gadus_morhua.gadMor3.0.dna.toplevel.csv};
		\addplot+[only marks, mark=*, mark size=1.5pt] table[col sep=comma, x=frequency, y=F_M_3]
		{./data/F_M_Gadus_morhua.gadMor3.0.dna.toplevel.csv};
		
		% Third panel
		\nextgroupplot[title={Human (raw)}]
		\addplot+[only marks, mark=*, mark size=1.5pt] table[col sep=comma, x=frequency, y=F_M_1]
		{./data/F_M_Homo_sapiens.GRCh38.dna.toplevel.csv};
		\addplot+[only marks, mark=*, mark size=1.5pt] table[col sep=comma, x=frequency, y=F_M_2]
		{./data/F_M_Homo_sapiens.GRCh38.dna.toplevel.csv};
		\addplot+[only marks, mark=*, mark size=1.5pt] table[col sep=comma, x=frequency, y=F_M_3]
		{./data/F_M_Homo_sapiens.GRCh38.dna.toplevel.csv};
		
	\end{groupplot}
	
	\node at (group c2r1.south) [below=1.2cm] {\pgfplotslegendfromname{sharedlegend}};
\end{tikzpicture}

	\caption{The percentage $F_M$ of minimizers with frequency up to $f$ on each level for the raw genomes Kestrel, Cod and Human. The percentages were measured using minimizers of length $17$ and the threshold $32$ for every level. Recall that the percentages $F_M$ are defined in relation to the number of $k$-mers left over on each level - not to the total number of $k$-mers (\cref{sec:minimizerDistribution_severalLevels}).} \label{fig:multilevelDistribution}
\end{figure}

It seems that super-$k$-mers with \frequent minimizers on one level are more likely to also have \frequent minimizers on subsequent levels.
A possible explanation for this observation are repetitively regions of the sequence in which a subsequences shorter than $k$ repeats many times over. $K$-mers in such regions would be likely to share the same minimizer regardless of the hash function used to compute it, resulting in those minimizers being frequent on every level. This hypothesis would have to be verified by searching for such repeating substrings in those super-$k$-mers which are frequent on all levels.\\

Minimizers being significantly less likely to be \infrequent on higher levels have a significant impact on the effectiveness of RSHash's leveled approach because they weaken its core functionality of subsequently exploiting minimizers the skewed and sparse distribution of minimizers.

\subsection{Operation Times} \label{subsec:observationsOperationTimes}
Recall that times 
$\timeForOperationOnStructure{\contains}{R}$,
$\timeForOperationOnStructure{\select}{S}$,
$\timeForOperationOnStructure{\access}{O}$,
$\timeForOperationOnStructure{\compare}{T}$ and
$\timeForOperationOnStructure{check}{H}$
needed for a rank query on bitvector $R$, a select queries on bitvector $S$, accessing an offset in $O$, a comparison with the text $T$ at a random offset and checking whether a key exists in the hash set $H$ are part of the model input and considered to be constant in the proposed model.\\
\begin{figure}[h]
	\centering
	\begin{minipage}{0.48\textwidth}
		\centering
		\resizebox{\linewidth}{!}{%
			\begin{tikzpicture}
            \begin{axis}[
                boxplot/draw direction=y,
                xtick={1,2,3,4,5},
                xticklabels={{$\timeForOperationOnStructure{\contains}{R}$},{$\timeForOperationOnStructure{\select}{S}$},{$\timeForOperationOnStructure{\access}{O}$},{$\timeForOperationOnStructure{\compare}{T}$},{$\timeForOperationOnStructure{check}{H}$}},
                x tick label style={align=center},
                ylabel={Measured Time (ns)},
            ]
            \addplot+ [draw=black, boxplot prepared={draw position=1, lower whisker=28.3553, lower quartile=101.1695, median=132.173, upper quartile=195.5475, upper whisker=335.796}] coordinates {};
\addplot+ [draw=black, boxplot prepared={draw position=2, lower whisker=8.29739, lower quartile=11.576775000000001, median=13.82165, upper quartile=19.037574999999997, upper whisker=30.1998}] coordinates {};
\addplot+ [draw=black, boxplot prepared={draw position=3, lower whisker=0.61052, lower quartile=2.114865, median=3.878525, upper quartile=5.1838049999999996, upper whisker=9.77881}] coordinates {};
\addplot+ [draw=black, boxplot prepared={draw position=4, lower whisker=2.30586, lower quartile=4.35256, median=6.587795, upper quartile=18.558425, upper whisker=39.858}] coordinates {};
\addplot+ [draw=black, boxplot prepared={draw position=5, lower whisker=2.33915, lower quartile=16.5644, median=23.736800000000002, upper quartile=27.63425, upper whisker=44.0691}] coordinates {};
\end{axis}
            \end{tikzpicture}
            
		}
		%\caption*{}
	\end{minipage}\hfill
	\begin{minipage}{0.48\textwidth}
		\centering
		\resizebox{\linewidth}{!}{%
			\begin{tikzpicture}
            \begin{axis}[
                boxplot/draw direction=y,
                xtick={1,2,3,4,5},
                xticklabels={{$\timeForOperationOnStructure{\contains}{R^1}$},{$\timeForOperationOnStructure{\contains}{R^2}$},{$\timeForOperationOnStructure{\contains}{R^3}$}},
                x tick label style={align=center},
                ylabel={Measured Time (ns)},
            ]
            \addplot+ [boxplot prepared={draw position=1, lower whisker=138.882, lower quartile=191.25799999999998, median=215.03, upper quartile=244.313, upper whisker=320.824}] coordinates {};
\addplot+ [boxplot prepared={draw position=2, lower whisker=36.8516, lower quartile=101.95275000000001, median=119.861, upper quartile=145.71125, upper whisker=211.199}] coordinates {};
\addplot+ [boxplot prepared={draw position=3, lower whisker=35.8562, lower quartile=85.84190000000001, median=100.5095, upper quartile=119.30225, upper whisker=169.385}] coordinates {};
\end{axis}
            \end{tikzpicture}
            
		}
		%\caption*{}
	\end{minipage}
	\caption{Box plots showing the distribution of measured operation times for all considered level parameters across all benchmark datasets.
		The left plot shows the combined distribution of the measured times 
		$\timeForOperationOnStructure{\contains}{R}$,
		$\timeForOperationOnStructure{\select}{S}$,
		$\timeForOperationOnStructure{\access}{O}$,
		$\timeForOperationOnStructure{\compare}{T}$,
		$\timeForOperationOnStructure{check}{H}$
		across all levels.
		The right plot shows the distribution of the measured times 
		$\timeForOperationOnStructure{\contains}{R^1}$,
		$\timeForOperationOnStructure{\contains}{R^2}$,
		$\timeForOperationOnStructure{\contains}{R^3}$
		needed for a rank query on bitvector $R$ on level one, two or three.}
	\label{fig:measureOperationTimes}
\end{figure}

\Cref{fig:measureOperationTimes} shows that the distributions of the measured times for these operation do not only depend on the level parameters but also substantially vary between levels. The time $\timeForOperationOnStructure{\contains}{R}$ for a rank query shows the highest variance with the lower ($101.17ns$) and upper ($195.55ns$) quartile across all measured values being almost $100ns$ apart. Its median across all levels is approximately $132.17ns$ and significantly varies when separating the measured values by levels one ($215.03ns$), two ($119.86ns$) and three ($100.51ns$).\\
The distribution of the measured times for $\timeForOperationOnStructure{\select}{S}$ and
$\timeForOperationOnStructure{\access}{O}$ have significantly smaller variances but also show declining medians for rising levels. This dependency of the level can be explained by the choice of only considering non-decreasing level parameters. This means that the vectors $R$, $S$ and $O$ on higher levels are significantly smaller than those on lower ones, allowing for faster query times.\\
The times $\timeForOperationOnStructure{\compare}{T}$ to compare a $k$-mer with the text at an offset and $\timeForOperationOnStructure{check}{H}$ to check whether the hash set contains a $k$-mer also appear to depend on the level parameters.\\

We will assess the implications of these observations in \cref{sec:modelAccuracy}. The values used for the model input are the observed medians of each operation time across all genomes, levels and level parameters:
\begin{align}
	\timeForOperationOnStructure{\contains}{R^l} &= 132.17ns\\
	\timeForOperationOnStructure{\select}{S^l} &= 13.82ns  \\
	\timeForOperationOnStructure{\access}{O^l} &= 3.88ns  \\
	\timeForOperationOnStructure{\compare}{T} &= 6.59ns   \\
	\timeForOperationOnStructure{check}{H} &= 23.74ns
\end{align}

\section{Model Accuracy} \label{sec:modelAccuracy}
\newcommand{\NAMEcorrCoef}{r}
\newcommand{\NAMEmeanErrRaw}{\bar{\mathcal{E}}}
\newcommand{\NAMEmeanErrAbs}{\bar{|\mathcal{E}|}}
\newcommand{\NAMEmeanErrSqrd}{\bar{\mathcal{E}^2}}
\newcommand{\corrCoef}{\NAMEcorrCoef = cov(Pred,Obs)/(\sigma_{Pred} \sigma_{Obs})}
\newcommand{\meanErrRaw}{\NAMEmeanErrRaw = \mathbb{E}(Pred-Obs)}
\newcommand{\meanErrAbs}{\NAMEmeanErrAbs = \mathbb{E}(|Pred-Obs|)}
\newcommand{\meanErrSqrd}{\NAMEmeanErrSqrd = \mathbb{E}((Pred-Obs)^2)}
We asses the accuracy of the lookup time and space consumption modeled by the MINLP for each of the datasets considered. To do this we first compute the model input and measure the actual space consumption and lookup time for every non-decreasing combination of level parameters on three levels as outlined in \cref{sec:computationOfModelInput}.\\
We then use a modified version of the proposed MINLP which forces an exact combination of level parameters to compute the modeled space consumption and lookup time for each combination of the level parameters considered on each raw and preprocessed sequence. We also refer to the modeled values as $Pred$ (predicted) and to the measured values as $Obs$ (observed).\\
We evaluate the models accuracy by comparing the predicted and observed values of the modeled variables using several statistical measures. Firstly, we compute the (Pearson) correlation coefficient $\corrCoef$ for each variable and dataset. This coefficient is invariant under separate changes in location and scale, rendering it a useful measure for linear dependence. It takes a value in $[-1,1] \cap \RR$ with coefficients close to $\pm 1$ indicating a strong correlation of the predicted and measured values whilst a coefficient close to $0$ can be interpreted as little to no correlation.\\
Furthermore we compute the mean error $\meanErrRaw$ to gain insights whether variables are being over or underestimated, the mean absolute error $\meanErrAbs$ to avoid cancellation of positive and negative errors and the mean squared error $\meanErrSqrd$ to indicate substantial outliers which would not show in $\NAMEmeanErrRaw$ and $\NAMEmeanErrAbs$.
\Cref{tab:modelErrorOverview} shows the these measures for the predicted and observed values of the lookup time total bits per $k$-mer across all datasets.\\
(In the following we present a selection of the computed measures aiming to point out those values deemed most relevant. For completeness all computed measures can be found in \cref{chap:completeModelEvaluationTables}, \crefrange{fig:correlation_stats_overview_kestrel_raw}{fig:correlation_stats_overview_human_unitigs}.)
\begin{table}[h]
	\centering
	\begin{tabular}{lrrrr}
\toprule
\textbf{$\timeForOperationOnStructure{lookup}{}$} & & & & \\
\midrule
Dataset & $\frac{cov(Pred,Obs)}{\sigma_{Pred} \sigma_{Obs}}$ & $\mathbb{E}(Pred-Obs)$ & $\mathbb{E}(|Pred-Obs|)$ & $\mathbb{E}(Pred-Obs)^2$ \\
\midrule
Kestrel (Raw) & 0.4840 & $-1.406 \cdot 10^{3}$ & $1.406 \cdot 10^{3}$ & $1.977 \cdot 10^{6}$ \\
Kestrel (Unitigs) & 0.3947 & $-1.698 \cdot 10^{3}$ & $1.698 \cdot 10^{3}$ & $2.884 \cdot 10^{6}$ \\
Cod (Raw) & 0.1406 & $-1.204 \cdot 10^{3}$ & $1.204 \cdot 10^{3}$ & $1.449 \cdot 10^{6}$ \\
Cod (Unitigs) & -0.0210 & $-1.546 \cdot 10^{3}$ & $1.546 \cdot 10^{3}$ & $2.389 \cdot 10^{6}$ \\
Human (Raw) & 0.2828 & $-1.872 \cdot 10^{3}$ & $1.872 \cdot 10^{3}$ & $3.506 \cdot 10^{6}$ \\
Human (Unitigs) & 0.2914 & $-2.609 \cdot 10^{3}$ & $2.609 \cdot 10^{3}$ & $6.806 \cdot 10^{6}$ \\
\midrule
\midrule
\textbf{$bpk_{total}$} & & & & \\
\midrule
Dataset & $\frac{cov(Pred,Obs)}{\sigma_{Pred} \sigma_{Obs}}$ & $\mathbb{E}(Pred-Obs)$ & $\mathbb{E}(|Pred-Obs|)$ & $\mathbb{E}(Pred-Obs)^2$ \\
\midrule
Kestrel (Raw) & 0.9995 & $7.453 \cdot 10^{-2}$ & $7.511 \cdot 10^{-2}$ & $5.642 \cdot 10^{-3}$ \\
Kestrel (Unitigs) & 0.9999 & $-1.271 \cdot 10^{-1}$ & $1.271 \cdot 10^{-1}$ & $1.616 \cdot 10^{-2}$ \\
Cod (Raw) & 0.7437 & $7.677 \cdot 10^{0}$ & $7.677 \cdot 10^{0}$ & $5.894 \cdot 10^{1}$ \\
Cod (Unitigs) & 0.9973 & $-2.380 \cdot 10^{-1}$ & $2.380 \cdot 10^{-1}$ & $5.664 \cdot 10^{-2}$ \\
Human (Raw) & 0.8530 & $8.637 \cdot 10^{0}$ & $8.637 \cdot 10^{0}$ & $7.460 \cdot 10^{1}$ \\
Human (Unitigs) & 0.9845 & $-2.092 \cdot 10^{-1}$ & $2.092 \cdot 10^{-1}$ & $4.377 \cdot 10^{-2}$ \\
\bottomrule
\end{tabular}
	\caption{Statistical evaluation of the variables modeling total bits per $k$-mer $bpk_{total}$ and lookup time $\timeForOperationOnStructure{lookup}{}$. The tables show the predicted (modeled) values $Pred$ compared against the observed (measured) ones $Obs$. The metrics used for comparison are the Pearson correlation coefficient the mean error, the mean absolute error and the mean squared error for each of the two variables. All values have been rounded to four significant figures.} 
	\label{tab:modelErrorOverview}
\end{table}

\subsection{Lookup Time}
The correlation coefficients of the lookup times $\timeForOperationOnStructure{lookup}{}$ predicted and observed values show that they correlate only weakly for the Kestrels genome (Raw: $\NAMEcorrCoef \approx 0.4840$, Unitigs: $\NAMEcorrCoef \approx 0.3947$) and little to none with for the Cods (Raw: $\NAMEcorrCoef \approx 0.1406$, Unitigs: $\NAMEcorrCoef \approx -0.0210$) and Humans (Raw: $\NAMEcorrCoef \approx 0.2828$, Unitigs: $\NAMEcorrCoef \approx 0.2914$) genome. 
Across all genomes we observe parameter combinations with almost identical predicted lookup times whose measured lookup times vary by several orders of magnitude. \Cref{fig:lookupTimePlot} shows this for the human genome.\\
\begin{figure}[h]
	\centering
	\begin{tikzpicture}
	\begin{groupplot}[
		group style={
			group size=2 by 1,
			horizontal sep=0.06\textwidth,
		},
		width=0.47\textwidth,
		height=5cm,
		]
		
		% First panel
		\nextgroupplot[
		title={Human (Raw)},
		ylabel={Observed $\timeForOperationOnStructure{lookup}{}$ (ns)},
		xlabel={Predicted $\timeForOperationOnStructure{lookup}{}$ (ns)},
		legend to name=sharedlegend,
		legend columns=7,
		, ymax=15000
		]
		\addplot[
		forget plot,
		only marks,
		scatter,
		point meta=explicit symbolic,
		scatter/classes={
			15={mark=*,mark size=1.5pt,draw=blue,fill=blue},
			17={mark=*,mark size=1.5pt,draw=red,fill=red},
			19={mark=*,mark size=1.5pt,draw=green!60!black,fill=green!60!black},
			21={mark=*,mark size=1.5pt,draw=orange,fill=orange},
			23={mark=*,mark size=1.5pt,draw=violet,fill=violet},
			25={mark=*,mark size=1.5pt,draw=cyan!70!black,fill=cyan!70!black}
		},
		] table[
		col sep=comma,
		x=predictedLookupTime,
		y=observedLookupTime,
		meta=m1,
		] {./data/sortedPredictedLookupTimes_Homo_sapiens.GRCh38.dna.toplevel.csv};
		\addlegendimage{empty legend}
		\addlegendentry{$\minimizerLength^1:$}
		
		\addlegendimage{only marks,mark=*,draw=blue,fill=blue}
		\addlegendentry{$15$}
		
		\addlegendimage{only marks,mark=*,draw=red,fill=red}
		\addlegendentry{$17$}
		
		\addlegendimage{only marks,mark=*,draw=green!60!black,fill=green!60!black}
		\addlegendentry{$19$}
		
		\addlegendimage{only marks,mark=*,draw=orange,fill=orange}
		\addlegendentry{$21$}
		
		\addlegendimage{only marks,mark=*,draw=violet,fill=violet}
		\addlegendentry{$23$}
		
		\addlegendimage{only marks,mark=*,draw=cyan!70!black,fill=cyan!70!black}
		\addlegendentry{$25$}
		

		
		% Second panel
		\nextgroupplot[title={Human (Unitigs)}, xlabel={Predicted $\timeForOperationOnStructure{lookup}{}$ (ns)}, ymax=15000]
		\addplot[
	    forget plot,
		only marks,
		scatter,
		point meta=explicit symbolic,
		scatter/classes={
			15={mark=*,mark size=1.5pt,draw=blue,fill=blue},
			17={mark=*,mark size=1.5pt,draw=red,fill=red},
			19={mark=*,mark size=1.5pt,draw=green!60!black,fill=green!60!black},
			21={mark=*,mark size=1.5pt,draw=orange,fill=orange},
			23={mark=*,mark size=1.5pt,draw=violet,fill=violet},
			25={mark=*,mark size=1.5pt,draw=cyan!70!black,fill=cyan!70!black}
		},
		] table[
		col sep=comma,
		x=predictedLookupTime,
		y=observedLookupTime,
		meta=m1,
		] {./data/sortedPredictedLookupTimes_Homo_sapiens.GRCh38.dna.toplevel.fa.unitigs.csv};

		
	\end{groupplot}
	\node at ($(group c1r1.south)!0.5!(group c2r1.south)$)
	[below=1.2cm] {\pgfplotslegendfromname{sharedlegend}};
\end{tikzpicture}
	\caption{Scatter plots comparing predicted lookup times $\timeForOperationOnStructure{Lookup}{}$ on the original scale with the observed values scaled by $10^4$ for each parameter combination on the raw (left) and preprocessed (right) human genomes. Each data point is labeled according to the minimizer length $\minimizerLength^1$ on level 1.} \label{fig:lookupTimePlot}
\end{figure}

At first it seems that the times $\timeForOperationOnStructure{\contains}{R}$,
$\timeForOperationOnStructure{\select}{S}$,
$\timeForOperationOnStructure{\access}{O}$,
$\timeForOperationOnStructure{\compare}{T}$ and
$\timeForOperationOnStructure{check}{H}$ which are constant parameters in the proposed MINLP actually vary significantly for different level parameters, levels and genomes as outlined in \cref{subsec:observationsOperationTimes} are a possible explanation for this. However, testing this hypothesis by computing the predicted lookup time using the measured operation times $\timeForOperationOnStructure{\contains}{R^\level}$,
$\timeForOperationOnStructure{\select}{S^\level}$,
$\timeForOperationOnStructure{\access}{O^\level}$ (for every level $1 \leq \level \leq 3$),
$\timeForOperationOnStructure{\compare}{T}$ and
$\timeForOperationOnStructure{check}{H}$ for every genome and parameter combination considered shows that this is not the case. \Cref{tab:WITH_MEASURED_OPERATION_TIMES_lookupTimeOverview} shows that computing the lookup times this way actually leads to even smaller correlation coefficients of the computed and observed values with mean errors similar to those when using constant operation times (\cref{tab:modelErrorOverview}).\\
\begin{table}[h]
	\centering
	\begin{tabular}{lrrrr}
\toprule
Dataset & $\NAMEcorrCoef$ & $\NAMEmeanErrRaw$ & $\NAMEmeanErrAbs$ & $\NAMEmeanErrSqrd$ \\
\midrule
Kestrel (Raw) & -0.3095 & $1.271 \times 10^{3}$ & $1.271 \times 10^{3}$ & $1.616 \times 10^{6}$ \\
Kestrel (Unitigs) & 0.0336 & $1.562 \times 10^{3}$ & $1.562 \times 10^{3}$ & $2.440 \times 10^{6}$ \\
Cod (Raw) & -0.1953 & $1.140 \times 10^{3}$ & $1.140 \times 10^{3}$ & $1.299 \times 10^{6}$ \\
Cod (Unitigs) & 0.0062 & $1.457 \times 10^{3}$ & $1.457 \times 10^{3}$ & $2.122 \times 10^{6}$ \\
Human (Raw) & -0.0231 & $1.695 \times 10^{3}$ & $1.695 \times 10^{3}$ & $2.873 \times 10^{6}$ \\
Human (Unitigs) & 0.0151 & $2.443 \times 10^{3}$ & $2.443 \times 10^{3}$ & $5.967 \times 10^{6}$ \\
\bottomrule
\end{tabular}
	\caption{Evaluation of the lookup times computed using the measured times $\timeForOperationOnStructure{\contains}{R^\level}$,
		$\timeForOperationOnStructure{\select}{S^\level}$,
		$\timeForOperationOnStructure{\access}{O^\level}$ (for every level $1 \leq \level \leq 3$),
		$\timeForOperationOnStructure{\compare}{T}$ and
		$\timeForOperationOnStructure{check}{H}$ instead of the proposed constant ones for every genome and parameter combination considered. The table shows the Pearson correlation coefficient, the mean error, the mean absolute error and the mean squared error computed for the predicted values compared against the observed ones. All values have been rounded to four significant figures.} 
	\label{tab:WITH_MEASURED_OPERATION_TIMES_lookupTimeOverview}
\end{table}

In addition to only correlating weakly, the observed lookup times are without exception three orders of magnitude larger than the predicted ones. This could possibly be due to components which count towards the lookup time but are not considered in the proposed model. One such example is the time needed to compute a $k$-mers minimizer; however, this alone does not explain the magnitude of the error.\\
Any constant offset to or linear scaling of either the modeled or measured times would not explain the observed inaccuracy regarding the lookup time as those would have no affect on the correlation coefficient which is invariant under separate changes in location and scale.\\

Understanding what causes the deviation between the modeled and measured lookup times to then account for it requires further investigation. If successful, this could considerably improve the models accuracy and therefore render the model more suitable for the optimization of level parameters.

\subsection{Space Consumption}
\Cref{tab:modelErrorOverview} shows that the smallest correlation coefficients across all considered sequences are $\NAMEcorrCoef \approx 0.7437$ for the Cods raw genome and $\NAMEcorrCoef \approx 0.8530$ for the raw Human genomes. These two genomes also show the most significant mean errors observed, namely $\NAMEmeanErrRaw \approx 7.677$ on the raw Cod and $\NAMEmeanErrRaw \approx 8.637$ on the raw Human genomes. Their respective mean absolute errors indicate that there are little to no negative errors for these values implying that the bits per $k$-mer for these two genomes are being strictly overestimated by approximately eight bits per $k$-mer.\\
The correlation coefficients regarding $bpk_{total}$ for the four remaining datasets are at least $\NAMEcorrCoef \approx 0.9845$ with mean absolute errors of at most $\NAMEmeanErrAbs \approx 0.2092$ bits per $k$-mer (both for Human (Unitigs)).
These observations can be explained by assessing, how well the proposed MINLP predicts each component of which the total bits per $k$-mer are composed, namely these are the bits per $k$-mer $bpk_{R^\level}$, $bpk_{S^\level}$ and $bpk_{O^\level}$ required for storing the vectors $R$, $S$ and $O$ on all considered levels as well as $bpk_T$ and $bpk_H$ needed to store the text and hash set.\\

For both the raw and preprocessed datasets the largest mean absolute error regarding the modeled and measured space consumption of the vectors $R^\level, S^\level$ and $O^\level$ on every level $1 \leq \level \leq 3$ are those of $bpk_{O^1}$, $bpk_{O^2}$ and $bpk_{O^3}$ for the preprocessed human genome. All three variables have a mean absolute error $\NAMEmeanErrAbs \approx 1.0161 \times 10^-2$. This shows that the proposed constraints for these variables model their space consumption accurately. However, it should be noted that this accuracy is only possible due to the precise computation of minimizer distributions.\\
These mean absolute errors of less than $0.02$ bits per $k$-mer are insignificant to the deviation of approximately eight bits per $k$-mer observed for the cods and humans raw genomes. Those must therefore stem from inaccuracies in the modeled $bpk_T$ and $bpk_H$.\\

\begin{table}[h]
	\centering
	\begin{tabular}{lrrrr}
\toprule
\textbf{$bpk_T$} & & & & \\
\midrule
Dataset & $\NAMEcorrCoef$ & $\NAMEmeanErrRaw$ & $\NAMEmeanErrAbs$ & $\NAMEmeanErrSqrd$ \\
\midrule
Kestrel (Raw) & -- & $6.350 \times 10^{-3}$ & $6.350 \times 10^{-3}$ & $4.032 \times 10^{-5}$ \\
Kestrel (Unitigs) & -- & $3.558 \times 10^{-2}$ & $3.558 \times 10^{-2}$ & $1.266 \times 10^{-3}$ \\
Cod (Raw) & -- & $1.300 \times 10^{-4}$ & $1.300 \times 10^{-4}$ & $1.690 \times 10^{-8}$ \\
Cod (Unitigs) & -- & $2.873 \times 10^{-1}$ & $2.873 \times 10^{-1}$ & $8.256 \times 10^{-2}$ \\
Human (Raw) & -- & $4.000 \times 10^{-5}$ & $4.000 \times 10^{-5}$ & $1.600 \times 10^{-9}$ \\
Human (Unitigs) & -- & $3.117 \times 10^{-1}$ & $3.117 \times 10^{-1}$ & $9.713 \times 10^{-2}$ \\
\midrule
\midrule
\textbf{$bpk_H$} & & & & \\
\midrule
Dataset & $\NAMEcorrCoef$ & $\NAMEmeanErrRaw$ & $\NAMEmeanErrAbs$ & $\NAMEmeanErrSqrd$ \\
\midrule
Kestrel (Raw) & 0.9051 & $-1.608 \times 10^{-1}$ & $1.608 \times 10^{-1}$ & $2.584 \times 10^{-2}$ \\
Kestrel (Unitigs) & 0.9767 & $1.178 \times 10^{-3}$ & $1.727 \times 10^{-3}$ & $2.986 \times 10^{-6}$ \\
Cod (Raw) & 0.6924 & $-7.729 \times 10^{0}$ & $7.729 \times 10^{0}$ & $5.974 \times 10^{1}$ \\
Cod (Unitigs) & 0.9642 & $2.706 \times 10^{-2}$ & $3.453 \times 10^{-2}$ & $1.192 \times 10^{-3}$ \\
Human (Raw) & 0.8991 & $-8.668 \times 10^{0}$ & $8.668 \times 10^{0}$ & $7.513 \times 10^{1}$ \\
Human (Unitigs) & 0.9694 & $3.890 \times 10^{-2}$ & $5.773 \times 10^{-2}$ & $3.333 \times 10^{-3}$ \\
\bottomrule
\end{tabular}
	\caption{Statistic evaluation for the variables modeling the total bits per $k$-mer $bpk_T$ for the text and $bpk_H$ for the hash set. The tables show the predicted values $Pred$ compared against the observed ones $Obs$. The metrics used for comparison are the Pearson correlation coefficient $\corrCoef$, the mean error, the mean absolute error and the mean squared error for each of the two variables. All values have been rounded to four significant figures. The correlation coefficient for the predicted and observed values of $bpk_T$ is not defined as the variance $\sigma_{Pred}$ of its predicted value $bpk_T = 2$ is zero.} 
	\label{tab:bpk_T_H_Overview}
\end{table}
\Cref{tab:bpk_T_H_Overview} shows that the observed errors mainly stem from an overestimation of the bits per $k$-mer $bpk_H$ needed to store the hash set. The overestimation $bpk_H$ by approximately $7.729$ (Cod) and $8.668$ (Human) bit for the cods and humans raw genomes are the main component of the large errors regarding their total space consumption $bpk_{total}$. The reason for this is that the raw sequences of both the human and the cod contain many duplicate $k$-mers. These can count towards the percentage of $k$-mers with \infrequent minimizers on all levels and thus lead to a larger value of the modeled $bpk_H$. However, in practice duplicates of the same $k$-mer collapse to a single entry in the hash set as the lookup only requires determining a single occurrence. This effect is negligent for the kestrels genome as it contains almost no duplicate $k$-mers.
\\
The comparatively small errors of the modeled $bpk_H$ for the preprocessed sequences is a result of the way the hash table is implemented in practice. Firstly, implementations of hash sets can contain metadata for each entry to enable faster queries at the cost of slightly higher space consumption. Secondly, hash sets are often initialized with a small number of buckets which doubles whenever a certain criteria regarding the hash sets load is reached. The resulting hash sets therefore have discrete sizes as can be seen for the cods genome in \cref{fig:bpkHCodSctterPlot}. These discrete sizes of hash sets with varying loads cannot be captured by a continues model which uses a constant load factor as proposed in \cref{constraint:bpkH}.\\
\begin{figure}[h]
	\centering
	\begin{tikzpicture}
	\begin{groupplot}[
		group style={
			group size=2 by 1,
			horizontal sep=0.06\textwidth,
		},
		width=0.47\textwidth,
		height=5cm,
		]
		
		% First panel
		\nextgroupplot[
		title={Cod (Raw)},
		ylabel={Observed $bpk_H$},
		xlabel={Predicted $bpk_H$},
		]
		\addplot[
		only marks,
		black,
		] table[
		col sep=comma,
		x=PredBpkH,
		y=ObsBpkH,
		] {./data/hashSetSizes_Gadus_morhua.gadMor3.0.dna.toplevel.csv};
		
		
		% Second panel
		\nextgroupplot[title={Cod (Unitigs)}, xlabel={Predicted $bpk_H$}]
		\addplot[
		only marks,
		black,
		] table[
		col sep=comma,
		x=PredBpkH,
		y=ObsBpkH,
		] {./data/hashSetSizes_Gadus_morhua.gadMor3.0.dna.toplevel.fa.unitigs.csv};
		
		
	\end{groupplot}
\end{tikzpicture}
	\caption{Scatter plots comparing predicted bits per $k$-mer $bpk_H$ of the hash set with the observed values for each parameter combination on the raw (left) and preprocessed (right) cod genomes. Each data point is labeled according to the minimizer length $\minimizerLength^1$ on level 1.} \label{fig:bpkHCodSctterPlot}
\end{figure}

\Cref{tab:bpk_T_H_Overview} shows that the errors regarding $bpk_T$ are significantly smaller than those of $bpk_H$. The largest observed errors for the bits per $k$-mer $bpk_T$ needed to store the text are those for the cods ($\NAMEmeanErrRaw \approx 0.2873$) and humans ($\NAMEmeanErrRaw \approx 0.3117$) preprocessed genome. This is again due to the repetitiveness of the respective raw genomes. The chosen approximation of $bpk_T \approx 2$ holds for the raw sequences because those are generally much longer than $k$. However, this is not necessarily the case when the input sequences have been preprocessed to remove duplicate $k$-mers. The preprocessed sequences are represented through unitigs, which are non-branching paths in the compact de Bruijn graph of the sequence \cite{doi:10.1089/cmb.2020.0431}. For highly repetitive genomes, their lengths are much closer to $k$, requiring more bits per $k$-mer to store. This effect is again negligent for the common kestrels genome which contains almost no duplicate $k$-mers. These errors could be eliminated by including the bits per $k$-mer needed to store a sequence in the MINLP input.

\section{Optimization} \label{sec:optimization}
\newcommand{\optGoal}[1]{optGoal_{#1}}
\newcommand{\minLookupTime}{\timeForOperationOnStructure{lookup}{}^{min}}
\newcommand{\minSpaceConsumption}{bpk_{total}^{min}}
To assess how well the proposed MINLP is suited for finding optimal level parameters we run a number of experiments. We specify four optimization goals for searching optimal parameters. Firstly, we search for level parameters minimizing the lookup time $\timeForOperationOnStructure{lookup}{}$ with no bound regarding the space consumption. Secondly, we search for level parameters minimizing the space consumption $bpk_{total}$ with no bound regarding the lookup time. Let $\minLookupTime$ be the minimal lookup time resulting from the first run and $\minSpaceConsumption$ be the minimal space consumption resulting from the second one. The minimal space consumption can then be used to specify the third optimization goal, minimizing $\timeForOperationOnStructure{lookup}{}$ whilst not exceeding the minimal space consumption $\minSpaceConsumption$ by more than fifty percent. Analogously the fourth optimization goal minimizes the space consumption without exceeding the minimal lookup time $\minLookupTime$ by more than fifty percent. \Cref{tab:optGoals} shows these four optimization goals.\\
\begin{table}[h]
	\centering
	\begin{tabular}{l|ccc}
		\toprule
		$\optGoal{}$ & $\minimize$ 	& $\boundTime$  & $\boundSpace$ \\
		\midrule
		$\optGoal{1}$ & $\timeForOperationOnStructure{lookup}{}$& $\infty$ & $\infty$ \\
		$\optGoal{2}$ & $bpk_{total}$ 							& $\infty$ & $\infty$ \\
		$\optGoal{3}$ & $\timeForOperationOnStructure{lookup}{}$& $\infty$ & $1.5 \cdot \minSpaceConsumption$ \\
		$\optGoal{4}$ & $bpk_{total}$ & $1.5 \cdot \minLookupTime$ & $\infty$ \\
		\bottomrule
	\end{tabular}
	\caption{The four considered optimization goals where $\minLookupTime$ is the minimal lookup time resulting from the first optimization and $\minSpaceConsumption$ is the minimal space consumption resulting from the second one.} \label{tab:optGoals}
\end{table}

We use the SCIP solver \cite{AchterbergBertholdKochetal.2008, Achterberg2009, BestuzhevaChmielaMuellerSerranoVigerskeWegscheider2023} which is capable of solving mixed integer non linear programs to solve the proposed MINLP for each specified sequence and optimization goal. This results
\todo[inline]{Tabelle gefundene Parameter}

\todo[inline]{Tabelle, Jonas gewählte Parameter}

\todo[inline]{Vergleich}


- number of levels variable instead of specified input



\chapter{Conclusion and Outlook} \label{chap:outlook}
\todo[inline]{erkenntnisse aus \cref{sec:optimization} einbringen}
The proposed mixed integer linear program does not satisfy the requirements for finding optimal level parameters to build RSHash for a given sequences. Whilst the MINLP accurately models the data structures space consumption, it is highly inaccurate regarding RSHash's lookup time. Further examination of these inaccuracies origins and possibly accounting for them could render the proposed MINLP more suitable to finding optimal level parameters.\\
A second significant restriction is that the proposed method requires specifying not only a number of levels but also a selection of level parameters to be considered. Furthermore, the introduced algorithm for computing the minimizer distributions required for the model input takes exponential time to compute, rendering its runtime only insignificantly faster than that of finding optimal level parameters by brute-forcing every considered combination of level parameters. Working towards a better understanding of the distribution of minimizers on several levels in dependence of level parameters could lead to faster methods for computing or approximating the required model input.\\
Even though not sufficient for finding optimal level parameters, attempting to model RSHash's lookup time and space consumption contributed towards a better understanding of the minimizer distributions influence on the data structures properties. One particularly interesting observations is how minimizer distribution on higher levels significantly depend on the level parameters of lower levels, weakening the repetitive compression using of minimizers.\\

In the case of resolving the described restrictions, there are several possibilities of broadening the proposed methods applicability.\\
Whilst the proposed MINLP models the time for so called positive lookups of $k$-mers which occur in the sequence, it could be extended to account for negative lookups of $k$-mers which do not appear in the sequence. This could be done by using a ratio of positive to negative lookups which positive lookups taking the proposed time and negative lookups consisting of a check in every levels $R$ vector followed by a check in the hash set.\\
Recall that whilst we examined RSHash build to answer lookup queries, the data structure can also be build to answer locate queries which return not only whether a $k$-mer exists but also all its offsets in the text. The proposed model could be extended to support modeling RSHash when build for answering locate queries. This would require modeling not the size of a hash set but that of a hash table, which includes a list of offsets for every $k$-mer that is frequent on all levels. The modeled lookup time could be adjusted to correspond to always searching through all offsets of a minimizer instead of only half of them on average.


\printbibliography[heading=bibintoc,title={Bibliography}]

\appendix
\chapter{Complete Model Evaluation Tables} \label{chap:completeModelEvaluationTables}
\Crefrange{fig:correlation_stats_overview_kestrel_raw}{fig:correlation_stats_overview_human_unitigs}
show the comparison of the modeled (predicted: $Pred$) values against the measured (observed: $Obs$) ones for every raw and preprocessed genome considered. Specifically each table shows the correlation coefficient $\corrCoef$, the mean error $\meanErrRaw$, the mean absolute error $\meanErrAbs$ and the mean squared error $\meanErrSqrd$ for each of the variables $bpk_{R^\level}$, $bpk_{S^\level}$ and $bpk_{O^\level}$ for all levels $1 \leq \level \leq 3$ as well as $bpk_T$, $bpk_H$, $bpk_{total}$ and $\timeForOperationOnStructure{lookup}{}$. Recall that the proposed MINLP (\cref{chap:minlp}) uses a constant approximation of $bpk_T = 2$. The correlation coefficient $\corrCoef$ is therefore not defined for the predicted and observed values of $bpk_T$ as the variance $\sigma_{Pred}$ of its predicted value is zero:
\begin{table} 
	\csvautotabular[separator=semicolon]{./data/correlation_stats_Falco_tinnunculus.FalTin1.0.dna.toplevel.csv}
	\caption{Kestrel (Raw)} \label{fig:correlation_stats_overview_kestrel_raw}
\end{table}

\begin{table} 
	\csvautotabular[separator=semicolon]{./data/correlation_stats_Falco_tinnunculus.FalTin1.0.dna.toplevel.fa.unitigs.csv}
	\caption{Kestrel (Unitigs)} \label{fig:correlation_stats_overview_kestrel_unitigs}
\end{table}

\begin{table} 
	\csvautotabular[separator=semicolon]{./data/correlation_stats_Gadus_morhua.gadMor3.0.dna.toplevel.csv}
	\caption{Cod (Raw)} \label{fig:correlation_stats_overview_cod_raw}
\end{table}

\begin{table} 
	\csvautotabular[separator=semicolon]{./data/correlation_stats_Gadus_morhua.gadMor3.0.dna.toplevel.fa.unitigs.csv}
	\caption{Cod (Unitigs)} \label{fig:correlation_stats_overview_cod_unitigs}
\end{table}


\begin{table} 
	\csvautotabular[separator=semicolon]{./data/correlation_stats_Homo_sapiens.GRCh38.dna.toplevel.csv}
	\caption{Human (Raw)} \label{fig:correlation_stats_overview_human_raw}
\end{table}

\begin{table} 
	\csvautotabular[separator=semicolon]{./data/correlation_stats_Homo_sapiens.GRCh38.dna.toplevel.fa.unitigs.csv}
	\caption{Human (Unitigs)} \label{fig:correlation_stats_overview_human_unitigs}
\end{table}

\end{document}
